% ===========================================================================
% handout-crypto.tex —— CTF 网络安全入门：Crypto 专攻版
% 重构依据：STYLE-SPEC.md（统一结构 / 题面卡 / 逐题带做 / 截图流水线 / 验收指标）
% 本册约定：只使用真实执行产生的截图与真实 flag；平台页面不伪造截图。
% ===========================================================================
\documentclass[UTF8,12pt,oneside]{ctexbook}
% handout-common.tex —— 五本讲义共享的导言区与排版宏
% 所有 handout-*.tex 通过 % handout-common.tex —— 五本讲义共享的导言区与排版宏
% 所有 handout-*.tex 通过 % handout-common.tex —— 五本讲义共享的导言区与排版宏
% 所有 handout-*.tex 通过 \input{handout-common} 引入本文件。
% 【重要】此文件为统一规范，任何单本讲义不得修改它。
% ===========================================================================
\usepackage[a4paper,left=18mm,right=18mm,top=18mm,bottom=18mm,
  includeheadfoot,headheight=15pt,headsep=4mm,footskip=10mm]{geometry}
\usepackage{amsmath,amssymb,booktabs,longtable,array,enumitem,listings,xcolor,hyperref,fancyhdr,graphicx}
\graphicspath{{../figures/}{../}}
\hypersetup{hidelinks,pdfauthor={CTF 培训资料}}

\setlength{\parindent}{2em}
\setlength{\parskip}{0.25em}
\emergencystretch=3em
\setlist{itemsep=0.15em,topsep=0.25em}
\lstset{basicstyle=\ttfamily\small,breaklines=true,columns=fullflexible,keepspaces=true,
  showstringspaces=false,frame=single,framerule=0.3pt,xleftmargin=0.4em,xrightmargin=0.4em}

% ---- 页眉页脚 -------------------------------------------------------------
\pagestyle{fancy}
\fancyhf{}
\fancyhead[L]{CTF 网络安全入门}
\fancyhead[R]{\thehandoutmark}
\fancyfoot[C]{\thepage}
\renewcommand{\headrulewidth}{0.3pt}
\newcommand{\thehandoutmark}{讲义}
% 用法（写在导言区）：\handoutmark{Web 专攻版}
\newcommand{\handoutmark}[1]{\renewcommand{\thehandoutmark}{#1}}
\newcommand{\booktitle}[1]{\hypersetup{pdftitle={#1}}}

% ---- 正文辅助宏 -----------------------------------------------------------
\newcommand{\goal}[1]{\noindent\textbf{本章目标：}#1\par}
\newcommand{\code}[1]{\texttt{#1}}
\newcommand{\kw}[1]{\textbf{#1}}
\newcommand{\flagbox}[1]{\par\noindent{\setlength{\fboxsep}{4pt}\fbox{\textbf{flag：}\texttt{#1}}}\par}
% 题面卡中的字段行：\field{附件}{...}
\newcommand{\field}[2]{\par\noindent\textbf{#1：}#2\par}

% ---- 信息框（一句话记忆 / 常见误区 / 排错指南 / 小技巧）--------------------
\newsavebox{\ibox}
\newenvironment{infobox}[1]
  {\par\medskip\begin{lrbox}{\ibox}\begin{minipage}{0.95\linewidth}\textbf{#1}\par\smallskip\hrule\smallskip}
  {\end{minipage}\end{lrbox}\noindent\fbox{\usebox{\ibox}}\par\medskip}
\newcommand{\memory}[1]{\begin{infobox}{一句话记忆}#1\end{infobox}}
\newcommand{\pitfall}[1]{\begin{infobox}{常见误区}#1\end{infobox}}
\newcommand{\debugtip}[1]{\begin{infobox}{排错指南}#1\end{infobox}}
\newcommand{\tip}[1]{\begin{infobox}{小技巧}#1\end{infobox}}

% ---- 题面卡 ---------------------------------------------------------------
% \begin{challenge}{题 R1：一句话标题} ... \field{...}{...} ... \end{challenge}
\newenvironment{challenge}[1]
  {\par\medskip\begin{lrbox}{\ibox}\begin{minipage}{0.95\linewidth}{\large\bfseries #1}\par\smallskip\hrule\smallskip}
  {\end{minipage}\end{lrbox}\noindent\fbox{\usebox{\ibox}}\par\medskip}

% ---- 插图 -----------------------------------------------------------------
% \shot[宽度比例]{文件名}{图注}{标签}   引用：\figref{标签}
\newcommand{\shot}[4][0.96]{%
  \begin{figure}[htbp]\centering
  \includegraphics[width=#1\linewidth]{#2}%
  \caption{#3}\label{fig:#4}\end{figure}}
\newcommand{\figref}[1]{图~\ref{fig:#1}}

\endinput
 引入本文件。
% 【重要】此文件为统一规范，任何单本讲义不得修改它。
% ===========================================================================
\usepackage[a4paper,left=18mm,right=18mm,top=18mm,bottom=18mm,
  includeheadfoot,headheight=15pt,headsep=4mm,footskip=10mm]{geometry}
\usepackage{amsmath,amssymb,booktabs,longtable,array,enumitem,listings,xcolor,hyperref,fancyhdr,graphicx}
\graphicspath{{../figures/}{../}}
\hypersetup{hidelinks,pdfauthor={CTF 培训资料}}

\setlength{\parindent}{2em}
\setlength{\parskip}{0.25em}
\emergencystretch=3em
\setlist{itemsep=0.15em,topsep=0.25em}
\lstset{basicstyle=\ttfamily\small,breaklines=true,columns=fullflexible,keepspaces=true,
  showstringspaces=false,frame=single,framerule=0.3pt,xleftmargin=0.4em,xrightmargin=0.4em}

% ---- 页眉页脚 -------------------------------------------------------------
\pagestyle{fancy}
\fancyhf{}
\fancyhead[L]{CTF 网络安全入门}
\fancyhead[R]{\thehandoutmark}
\fancyfoot[C]{\thepage}
\renewcommand{\headrulewidth}{0.3pt}
\newcommand{\thehandoutmark}{讲义}
% 用法（写在导言区）：\handoutmark{Web 专攻版}
\newcommand{\handoutmark}[1]{\renewcommand{\thehandoutmark}{#1}}
\newcommand{\booktitle}[1]{\hypersetup{pdftitle={#1}}}

% ---- 正文辅助宏 -----------------------------------------------------------
\newcommand{\goal}[1]{\noindent\textbf{本章目标：}#1\par}
\newcommand{\code}[1]{\texttt{#1}}
\newcommand{\kw}[1]{\textbf{#1}}
\newcommand{\flagbox}[1]{\par\noindent{\setlength{\fboxsep}{4pt}\fbox{\textbf{flag：}\texttt{#1}}}\par}
% 题面卡中的字段行：\field{附件}{...}
\newcommand{\field}[2]{\par\noindent\textbf{#1：}#2\par}

% ---- 信息框（一句话记忆 / 常见误区 / 排错指南 / 小技巧）--------------------
\newsavebox{\ibox}
\newenvironment{infobox}[1]
  {\par\medskip\begin{lrbox}{\ibox}\begin{minipage}{0.95\linewidth}\textbf{#1}\par\smallskip\hrule\smallskip}
  {\end{minipage}\end{lrbox}\noindent\fbox{\usebox{\ibox}}\par\medskip}
\newcommand{\memory}[1]{\begin{infobox}{一句话记忆}#1\end{infobox}}
\newcommand{\pitfall}[1]{\begin{infobox}{常见误区}#1\end{infobox}}
\newcommand{\debugtip}[1]{\begin{infobox}{排错指南}#1\end{infobox}}
\newcommand{\tip}[1]{\begin{infobox}{小技巧}#1\end{infobox}}

% ---- 题面卡 ---------------------------------------------------------------
% \begin{challenge}{题 R1：一句话标题} ... \field{...}{...} ... \end{challenge}
\newenvironment{challenge}[1]
  {\par\medskip\begin{lrbox}{\ibox}\begin{minipage}{0.95\linewidth}{\large\bfseries #1}\par\smallskip\hrule\smallskip}
  {\end{minipage}\end{lrbox}\noindent\fbox{\usebox{\ibox}}\par\medskip}

% ---- 插图 -----------------------------------------------------------------
% \shot[宽度比例]{文件名}{图注}{标签}   引用：\figref{标签}
\newcommand{\shot}[4][0.96]{%
  \begin{figure}[htbp]\centering
  \includegraphics[width=#1\linewidth]{#2}%
  \caption{#3}\label{fig:#4}\end{figure}}
\newcommand{\figref}[1]{图~\ref{fig:#1}}

\endinput
 引入本文件。
% 【重要】此文件为统一规范，任何单本讲义不得修改它。
% ===========================================================================
\usepackage[a4paper,left=18mm,right=18mm,top=18mm,bottom=18mm,
  includeheadfoot,headheight=15pt,headsep=4mm,footskip=10mm]{geometry}
\usepackage{amsmath,amssymb,booktabs,longtable,array,enumitem,listings,xcolor,hyperref,fancyhdr,graphicx}
\graphicspath{{../figures/}{../}}
\hypersetup{hidelinks,pdfauthor={CTF 培训资料}}

\setlength{\parindent}{2em}
\setlength{\parskip}{0.25em}
\emergencystretch=3em
\setlist{itemsep=0.15em,topsep=0.25em}
\lstset{basicstyle=\ttfamily\small,breaklines=true,columns=fullflexible,keepspaces=true,
  showstringspaces=false,frame=single,framerule=0.3pt,xleftmargin=0.4em,xrightmargin=0.4em}

% ---- 页眉页脚 -------------------------------------------------------------
\pagestyle{fancy}
\fancyhf{}
\fancyhead[L]{CTF 网络安全入门}
\fancyhead[R]{\thehandoutmark}
\fancyfoot[C]{\thepage}
\renewcommand{\headrulewidth}{0.3pt}
\newcommand{\thehandoutmark}{讲义}
% 用法（写在导言区）：\handoutmark{Web 专攻版}
\newcommand{\handoutmark}[1]{\renewcommand{\thehandoutmark}{#1}}
\newcommand{\booktitle}[1]{\hypersetup{pdftitle={#1}}}

% ---- 正文辅助宏 -----------------------------------------------------------
\newcommand{\goal}[1]{\noindent\textbf{本章目标：}#1\par}
\newcommand{\code}[1]{\texttt{#1}}
\newcommand{\kw}[1]{\textbf{#1}}
\newcommand{\flagbox}[1]{\par\noindent{\setlength{\fboxsep}{4pt}\fbox{\textbf{flag：}\texttt{#1}}}\par}
% 题面卡中的字段行：\field{附件}{...}
\newcommand{\field}[2]{\par\noindent\textbf{#1：}#2\par}

% ---- 信息框（一句话记忆 / 常见误区 / 排错指南 / 小技巧）--------------------
\newsavebox{\ibox}
\newenvironment{infobox}[1]
  {\par\medskip\begin{lrbox}{\ibox}\begin{minipage}{0.95\linewidth}\textbf{#1}\par\smallskip\hrule\smallskip}
  {\end{minipage}\end{lrbox}\noindent\fbox{\usebox{\ibox}}\par\medskip}
\newcommand{\memory}[1]{\begin{infobox}{一句话记忆}#1\end{infobox}}
\newcommand{\pitfall}[1]{\begin{infobox}{常见误区}#1\end{infobox}}
\newcommand{\debugtip}[1]{\begin{infobox}{排错指南}#1\end{infobox}}
\newcommand{\tip}[1]{\begin{infobox}{小技巧}#1\end{infobox}}

% ---- 题面卡 ---------------------------------------------------------------
% \begin{challenge}{题 R1：一句话标题} ... \field{...}{...} ... \end{challenge}
\newenvironment{challenge}[1]
  {\par\medskip\begin{lrbox}{\ibox}\begin{minipage}{0.95\linewidth}{\large\bfseries #1}\par\smallskip\hrule\smallskip}
  {\end{minipage}\end{lrbox}\noindent\fbox{\usebox{\ibox}}\par\medskip}

% ---- 插图 -----------------------------------------------------------------
% \shot[宽度比例]{文件名}{图注}{标签}   引用：\figref{标签}
\newcommand{\shot}[4][0.96]{%
  \begin{figure}[htbp]\centering
  \includegraphics[width=#1\linewidth]{#2}%
  \caption{#3}\label{fig:#4}\end{figure}}
\newcommand{\figref}[1]{图~\ref{fig:#1}}

\endinput

% Keep instructional figures beside the step that introduces them.
\usepackage{float}
\renewcommand{\shot}[4][0.92]{%
  \begin{figure}[H]\centering
  \includegraphics[width=#1\linewidth]{#2}%
  \caption{#3}\label{fig:#4}\end{figure}}
% 字形修正：①②③④ 等序号、★☆ 与希腊字母在西文字体中缺字形，声明为 CJK 字符类
% 改由中文字体渲染（↔ 中西文字体都缺，用 $\leftrightarrow$ 排版）。
\xeCJKDeclareCharClass{CJK}{"2460 -> "24FF, "2194 -> "2195, "2605 -> "2606, "03B1 -> "03C9}
\handoutmark{Crypto 专攻版}
\booktitle{CTF 网络安全入门：Crypto专攻版}

\begin{document}
\frontmatter

\begin{titlepage}
\centering
\vspace*{32mm}
{\Huge\bfseries CTF 网络安全入门：Crypto 专攻版\par}
\vspace{9mm}
{\LARGE 先识别表示，再找可逆关系\par}
\vspace{5mm}
{\large 零基础版 · 题目 C0--C4（含玄机 586 五格电文、587 联号回执）\par}
\vspace{14mm}
{\large 本册特色\par}
\vspace{3mm}
\begin{minipage}{0.86\linewidth}
\begin{itemize}
\item 零基础预备篇：字节、十六进制、异或、模运算、RSA 三元组，每个术语先讲人话；
\item 每题一张标准题面卡：给什么、要什么、交付物、考点、前置知识、难度；
\item 逐题图解带做：每一步命令配真实终端截图，看到图就能照着做；
\item 所有截图来自真实执行（termcap 实录 + 终端渲染），所有 flag 来自真实运行输出；
\item 附命令速查表、术语表、自测题与完整推演。
\end{itemize}
\end{minipage}
\vfill
{\large 适用对象：学过 Python 和命令行、第一次接触 CTF 的同学\par}
\vspace{4mm}
{\large 版本：2026 年 9 月\par}
\end{titlepage}

\chapter*{使用说明}
\addcontentsline{toc}{chapter}{使用说明}

\section*{这本讲义带你做什么}
Crypto（密码学）方向的题目，说到底只有两种任务：\kw{把看不懂的数据还原成看得懂的数据}，或者\kw{把一段关系想明白}。本册用 5 道题把最常见的三类突破口各练一遍：数据表示（C0）、异或与穷举（C1）、RSA 的数学关系（C2--C4）。做完之后，你应当能自己回答三个问题：手里这份数据到底是什么表示？它的密钥空间有多大？题目给了什么可验证的关系？

\section*{学习路线图}
\begin{enumerate}
\item \kw{第 1 章 零基础预备}：装好 Python 3 和 xxd，跑通第一条命令，确认你能看到和讲义一样的输出；
\item \kw{第 2 章 Crypto 基础知识铺垫}：字节与十六进制、异或真值表、模运算、RSA 三元组 (n,e,c)、小模数与共模的危险、自定义编码的方法论。这章不写题，但每一小节都对应后面某道题的某一步；
\item \kw{第 3 章 Crypto 工具箱}：python 一行命令、xxd、\code{pow} 三参数快速幂与求逆元，每个工具配真实输出截图；
\item \kw{第 4 章 题目：题面卡与逐题图解带做}：每题一张题面卡，卡后面直接就是这道题的解题步骤，按“读题与思路 → 分步操作与观察 → 结果与核对 → 为什么会成功 → 排错指南 → 变式练习”推进。建议先遮住解题步骤自己做一遍，卡住再看；
\item \kw{第 6、7 章}：复盘知识地图与易错清单，做自测题并对照完整推演；
\item \kw{附录}：玄机平台真题状态与提交 checklist、命令速查表、术语表。
\end{enumerate}

\section*{怎么用这本讲义}
\begin{itemize}
\item \kw{先看图再敲命令}：每张终端截图都标了这一步在做什么。截图里的路径以本资料根目录为当前目录，例如 \code{labs/crypto/xor.hex}。
\item \kw{每步都问“为什么”}：正文里每条命令后面都紧跟一句“这一步是为了……”。能自己说出为什么，才算会做。
\item \kw{核对不能省}：解出 flag 后一定要做“结果与核对”里的独立验证。脚本打印出 \code{flag\{...\}} 不等于答案正确。
\item \kw{变式练习是升级关}：做对原题只是及格，做完变式才算掌握。
\end{itemize}

\section*{安全声明}
所有实验只针对本资料附带的程序和明确允许练习的平台；不要把命令用于未经授权的系统。玄机平台入口为 \url{https://xj.edisec.net/}，授课时由教师提供个人或课程账号，讲义不保存账号口令。平台题面、价格与界面可能变化，始终以当前题目详情页为准。

\section*{截图说明}
本册所有终端截图都来自真实执行：先用 \code{tools/termcap.py} 在真实 shell 里执行命令并记录“命令 + 真实输出”，再用 \code{tools/term2png.py} 渲染成终端窗口图。每张图对应的原始实录保存在 \code{figures/raw/} 下的同名 \code{.txt} 文件里，可自行核对。原理示意图（异或真值表、RSA 流程图、五格电文比特流图）由 matplotlib 绘制，画的是结构关系。玄机平台的网页没有伪造截图，涉及平台状态的地方用文字表格说明。

\section*{术语约定}
每个术语第一次出现时都会用白话解释，之后直接使用。最容易混淆的三个词先放在这里：
\begin{itemize}
\item \kw{字符}：人眼看到的一个字，如字母 \code{A}、汉字“夜”；
\item \kw{字节}：计算机存储的一个 0--255 的数，是数据的最小常用单位；
\item \kw{十六进制}：字节的一种显示方式，一个字节显示成两个字符，如字节 65 显示成 \code{41}。
\end{itemize}

\tableofcontents
\mainmatter

% ===========================================================================
\chapter{零基础预备：实验环境与第一次动手}
% ===========================================================================
\goal{装好 Python 3 与 xxd；能在正确目录下运行命令；理解“文件内容”和“文件显示”不是一回事。}

\section{Crypto 题长什么样}
Crypto 题的题面通常很短：给一段“看不懂的数据”，要一个 \code{flag\{...\}} 字符串。看不懂的数据可能是一串十六进制字符、一堆整数、一串奇怪符号，或者几个大整数。破解它们不靠猜，靠三步：
\begin{enumerate}
\item \kw{识别表示}：这份数据现在是文本、字节、整数列表还是符号串？
\item \kw{找可逆关系}：加密过程有没有留下结构？密钥空间小不小？有没有重复使用参数？
\item \kw{验证}：还原结果能不能被独立证据确认（已知样本、再加密比对、平台反馈）？
\end{enumerate}
本册的 5 道题分别训练这三步里的不同环节。请记住第一句话：\kw{先识别表示，再找可逆关系}。

\section{安装与验证（Windows 侧）}
\paragraph{要装什么}
\begin{itemize}
\item \kw{Python 3}（3.8 以上即可）：所有解题脚本都用它。安装时勾选“Add python.exe to PATH”。
\item \kw{PowerShell}：Windows 自带，不用装。
\item \kw{WSL/Ubuntu}（可选，做 xxd 练习时用）：Microsoft Store 安装 Ubuntu-24.04。没有 WSL 也能做完本册全部题目，xxd 只是看十六进制的工具之一。
\end{itemize}

\paragraph{验证命令}在 PowerShell 里执行（当前目录切到本资料根目录）：
\begin{lstlisting}
python --version
python -c "print('hello, Crypto!'); print(2+3)"
Get-ChildItem labs/crypto | Select-Object Name, Length
\end{lstlisting}
应该看到 Python 3.x 的版本号、两行打印，以及 \code{labs/crypto} 目录下的三个文件：\code{make\_challenges.py}、\code{rsa.txt}、\code{xor.hex}。\figref{env-win} 是一次真实执行的截图。

\shot[0.95]{fig-crypto-01-env-check.png}{第 1 章 Windows 侧环境验证：python 版本、一行命令试跑、实验目录文件列表（真实终端实录）}{env-win}

\memory{看到和讲义一样的输出，才说明环境对了；输出不一样时先别做题，先修环境。}

\section{安装与验证（WSL 侧，可选）}
WSL 里的 Ubuntu 提供 \code{xxd}（十六进制查看工具）和 \code{python3}。验证命令（在 PowerShell 里执行）：
\begin{lstlisting}
wsl -d Ubuntu-24.04 -- bash -lc "python3 --version"
wsl -d Ubuntu-24.04 -- bash -lc "xxd -v"
\end{lstlisting}
\figref{env-wsl} 是真实执行结果：Ubuntu 里是 Python 3.12.3，\code{xxd} 为 2023-10-25 版。注意 Windows 的 \code{C:\textbackslash{}} 盘在 WSL 中通常位于 \code{/mnt/c/}，比如本资料根目录在 WSL 里是 \code{/mnt/c/Users/你的用户名/Desktop/CTF零基础自学资料包}。

\shot[0.95]{fig-crypto-02-env-check-wsl.png}{第 1 章 WSL 侧环境验证：python3 与 xxd 版本、实验目录列表（真实终端实录）}{env-wsl}

\pitfall{Windows 的 \code{python} 和 WSL 的 \code{python3} 是两套环境。本册命令若没有特别说明，都在 Windows PowerShell 里执行；标注“WSL”的才去 WSL。}

\section{目录与文件约定}
书中的相对路径以本资料根目录为当前目录。终端的当前目录不对时，程序会报“找不到文件”，这不代表题目坏了。先用 \code{pwd}（PowerShell 里是 \code{Get-Location}）和 \code{ls}（PowerShell 里是 \code{Get-ChildItem}）确认位置。本册用到的文件：
\begin{itemize}
\item \code{labs/crypto/xor.hex}：C1 的密文（十六进制文本文件）；
\item \code{labs/crypto/rsa.txt}：C3 的小模数 RSA 参数与密文列表；
\item \code{labs/crypto/make\_challenges.py}：生成上面两个文件的脚本，可用于复现题目数据（默认不要改动 \code{labs/} 下任何文件）；
\item \url{labs/platform/attachments/586-five-grid.zip}、\url{labs/platform/attachments/587-linked-receipts.zip}：两道平台题的官方附件本地副本；
\item \code{labs/platform/solve\_five\_grid.py}、\code{labs/platform/solve\_receipt.py}：两道平台题的复现脚本。
\end{itemize}

\section{第一次运行脚本}
在根目录运行平台题脚本（现在只是“试跑”，第 4 章会讲每一步的含义）：
\begin{lstlisting}
python labs/platform/solve_five_grid.py
python labs/platform/solve_receipt.py
\end{lstlisting}
如果两条命令都打印出 \code{flag: flag\{...\}} 字样，说明 Python 环境、路径、附件副本全部就绪。看到 flag 不用记，后面会带你一步步把它们做出来。

\memory{脚本能跑通 = 环境就绪；脚本报“找不到文件” = 当前目录不对；脚本报编码错误 = 先加 \code{python -X utf8} 再试。}

\section{用独立证据验证答案}
脚本打印一串 \code{flag\{...\}} 只说明脚本输出了字符串，不等于答案正确。本册每个“结果与核对”模块都给出\kw{独立于最后那行打印}的验证方法：
\begin{itemize}
\item C1：把明文重新异或一遍，看是否还原出原密文；
\item C3：用公钥把明文再加密一遍，看是否等于密文列表；
\item C2（586）：先用公开试报解出已知明文“夜班OK”，映射对了才解正文；
\item C4（587）：用两个指数分别再加密，两份都要对上；
\item 平台题最后还要看提交页面是否显示“FLAG 正确”。
\end{itemize}
写题解时也应写出从附件条件到结果的证据链，而非只附一行看似成功的输出。

% ===========================================================================
\chapter{Crypto 基础知识铺垫}
% ===========================================================================
\goal{把字节、十六进制、异或、模运算、RSA 三元组、小模数与共模、自定义编码这七块地基打牢，每块都能对应到后面的具体题目。}

\section{计算机里的数据：字符、字节、十六进制}
\kw{字符}是人看到的字；\kw{字节}是计算机存的一个 0--255 的数；\kw{十六进制}是字节的一种显示方式。三者的关系是：字符 $\xrightarrow{\text{编码}}$ 字节 $\xrightarrow{\text{显示}}$ 十六进制文本。

\paragraph{ASCII}最古老的编码表：把 0--127 的数分配给英文字母、数字、标点。例如字母 \code{A} 对应 65（十六进制 \code{41}），小写 \code{f} 对应 102（\code{66}）。英文字符在 ASCII 里占 1 个字节。\kw{UTF-8}是目前最常用的编码：英文部分兼容 ASCII（1 字节），汉字通常占 3 字节。

\paragraph{动手看一眼}运行 \figref{ascii-hex} 里的两条命令：第一条把字符 \code{A} 编码成字节再显示成十六进制；第二条把 \code{flag\{} 这五个字符逐个打印成十六进制。你应该看到：\code{b'A' 41 1} 表示字符 \code{A} 是 1 个字节、十六进制写作 \code{41}；\code{b'41'} 却是 2 个字节——那是字符 \code{4} 和字符 \code{1}！而 \code{bytes.fromhex('41')} 才是 1 个字节。

\shot[0.95]{fig-crypto-03-ascii-hex-demo.png}{第 2 章 字节、ASCII 与十六进制的关系：字符$\leftrightarrow$字节$\leftrightarrow$十六进制的换算演示（真实终端实录）}{ascii-hex}

\paragraph{最常踩的坑}把“十六进制文本”当成“字节本身”。文件里存着 44 个字符 \code{515b...544a}，它其实只是一段文本；用 \code{bytes.fromhex()} 解析后才得到 22 个字节。每跨过一种表示，都应写下转换方法和转换后的长度。

\pitfall{\code{b'41'} 是 2 个字节（两个 ASCII 字符），\code{bytes.fromhex('41')} 是 1 个字节（值为 0x41）。数字符个数前先想清楚手里的是哪种。}

\paragraph{表示方式的候选假设}拿到一串看不懂的数据，先做假设再验证：
\begin{itemize}
\item 只包含 \code{0-9a-fA-F} 且长度为偶数：可能是字节的十六进制表示；
\item 包含大小写字母、数字、斜线和加号，末尾有等号：可能是 Base64；
\item 固定长度的长十六进制串（比如 32 位）：可能是哈希摘要；
\item 由 32 个互不相同的怪符号组成：可能是自定义的 5 位编码（见 2.8 节）。
\end{itemize}
它们只是候选假设，不是确证。先试解析，再检查结果里有没有文件头、可打印文本或下一层编码。Base64 的转换规则公开、没有秘密密钥，所以它是\kw{编码}，不是\kw{加密}。

\paragraph{十六进制和 Base64 到底差在哪}两者都是“把字节换成一串字符”的\kw{编码}，规则公开、没有密钥，都能无损还原；区别只在\kw{用哪些字符}和\kw{膨胀多少}：
\begin{itemize}
\item \kw{十六进制}：1 个字节写成 2 个十六进制字符（长度变为 2 倍），字符表只有 \code{0-9a-f}；
\item \kw{Base64}：每 3 个字节写成 4 个可打印字符（长度约 1.33 倍），会出现大写、小写、数字和 \code{+} \code{/}，末尾常用 \code{=} 补齐。
\end{itemize}
识别时先看\kw{字符表}，再看\kw{长度比}。\figref{b64-hex} 用同一份 9 字节数据演示：\code{.hex()} 得到 18 个字符，\code{base64.b64encode()} 得到 12 个字符，二者都能还原原始字节。以后在 C0、C1 里看到 \code{0-9a-f} 就想到十六进制，看到 \code{+} \code{/} \code{=} 就想到 Base64。

\shot[0.95]{fig-crypto-25-b64-vs-hex.png}{第 2 章 Base64 与十六进制对照：同一份字节的两种显示、字符表与长度比的差别（真实终端实录）}{b64-hex}

\pitfall{同一份字节，十六进制和 Base64 的字面完全不同（例如 \code{666c6167} 与 \code{ZmxhZw==} 是同一段内容），别把两者混用：用错解码函数会报错或解出乱码。判断“是不是 Base64”不能只看长度，关键看字符表里有没有 \code{+} \code{/} \code{=}。}

\memory{字符是给人看的，字节是给计算机存的，十六进制和 Base64 都是给眼睛看字节的方式：前者 2 倍长、只有 \code{0-9a-f}，后者约 1.33 倍长、含 \code{+} \code{/} \code{=}。跨表示先记长度。}

\section{异或 XOR：真值表与“加密 = 解密”}
\kw{异或}（XOR，符号 $\oplus$）是一个按位运算：两个位相同得 0，不同得 1。它的真值表和可逆性见 \figref{xor-truth}。

\shot[0.95]{fig-crypto-04-xor-truth.png}{第 2 章 异或真值表与可逆性：两位比较规则，以及同一密钥做两次即还原（原理示意图）}{xor-truth}

为什么说“加密等于解密”？因为异或有个关键性质：
\[
(x \oplus k) \oplus k = x
\]
用密钥 $k$ 异或一次是加密，再用同一个 $k$ 异或一次就还原。所以异或加密的解密程序就是加密程序本身。Python 里逐字节异或写作 \code{bytes(x \textasciicircum{} k for x in data)}（\code{\textasciicircum} 是 Python 的异或运算符）。

\paragraph{异或的局限}若密钥与消息等长、完全随机、只使用一次（一次性密码本），理论上不可破。但 CTF 题里密钥通常很短、反复使用：C1 就是 1 个字节的密钥循环作用于所有明文字节。此时密文的结构就泄漏了密钥的结构。

\memory{异或一次是加密，异或两次是还原；密钥重复使用，结构就会泄漏。}

\section{为什么单字节异或可以穷举}
密钥只有 1 个字节，取值范围 0--255，一共只有 \kw{256 种可能}。哪怕不懂数学，把 256 种全试一遍也就试完了，这叫\kw{穷举}（brute force）。难点只在于：试出来的 256 个结果里，哪一个是真明文？判断方法叫\kw{可读性判断}：
\begin{itemize}
\item \kw{看前缀}：flag 通常以 \code{flag\{} 开头；
\item \kw{看字符范围}：真明文多半是可打印字符（ASCII 32--126），乱码会混进控制字符；
\item \kw{看语言}：英文明文里字母、空格占比高；中文明文按 UTF-8 解码要合法；
\item \kw{看结构}：花括号配对、单词可读、长度合理。
\end{itemize}
把候选密钥、可读比例和解出的原始字节都打印出来，方便复核。第 4 章 C1 会完整走一遍。

\pitfall{短数据可能有很多“看起来可读”的结果。只有一个字符被猜中不算数，至少要检查完整的 \code{flag\{} 前缀和整段的字符范围。}

\section{模运算：余数的世界}
\kw{模运算}（mod）就是“除以 $n$ 取余数”：$17 \bmod 5 = 2$。它有几个好用的性质：
\begin{itemize}
\item 结果永远落在 $0$ 到 $n-1$ 之间，再大的数取模后都变小；
\item 加、乘、乘方都可以先算再取模，结果不变。这就是\kw{快速幂}：算 $2^{100} \bmod 1000$ 不用先算出 31 位的 $2^{100}$，Python 的 \code{pow(2, 100, 1000)} 三参数版本专门干这个；
\item 在模 $n$ 意义下，有的数有\kw{逆元}：满足 $a \cdot a^{-1} \equiv 1 \pmod n$ 的那个 $a^{-1}$。Python 直接写 \code{pow(a, -1, n)}。
\end{itemize}
\paragraph{用钟表理解“绕回来”}模运算可以想成看钟面：一圈 12 小时，15 点其实就是下午 3 点，写成 $15 \bmod 12 = 3$；再走满一圈又回到 3 点——这正是“结果只落在 $0$ 到 $n-1$，一超过就从 0 重新数”。负数也一样：$(-8) \bmod 5 = 2$，因为从 0 往回退 8 步，绕了 2 圈之后落回 2。有了这个直觉就能明白：RSA 里的指数无论多大，取模后都会被“绕”回一个小区间，所以快速幂能一直保持小数字、算得又快又准。
\figref{mod-demo} 是真实演示：注意 $-8 \bmod 5 = 2$（Python 的余数永远非负），以及 $2$ 的幂次模 $5$ 会循环出现 $1,2,4,3$。求逆元并不总是可行：$a$ 与 $n$ 不互素时逆元不存在。

\paragraph{最大公约数与“互素”}刚刚反复出现的“互素”需要先定义。\kw{最大公约数} $\gcd(a, b)$ 指能同时整除 $a$ 和 $b$ 的最大正整数，Python 写作 \code{math.gcd(a, b)}，例如 \code{math.gcd(12, 18) = 6}。当 $\gcd(a, b) = 1$（两数除 1 以外没有公共因子）时，称 $a$ 与 $b$ \kw{互素}。互素是逆元存在的唯一前提：只有 $\gcd(a, n) = 1$，\code{pow(a, -1, n)} 才求得出逆元，否则会报 \code{ValueError}。“$e$ 与 $\varphi(n)$ 互素”这道 RSA 门槛，用的就是这个概念。

\shot[0.95]{fig-crypto-05-mod-demo.png}{第 2 章 模运算演示：余数、快速幂、循环性（真实终端实录）}{mod-demo}

\memory{模 $n$ = 除以 $n$ 取余数；快速幂用 \code{pow(b, e, n)}；求逆元用 \code{pow(a, -1, n)}，前提是 $a$ 与 $n$ 互素。}

\section{RSA 从头讲：三元组 $(n, e, c)$ 与私钥 $d$}
RSA 是最常见的公钥加密算法。CTF 题里你通常会拿到三个数：\kw{模数 $n$}、\kw{加密指数 $e$}、\kw{密文 $c$}，合称三元组 $(n, e, c)$。加密与解密流程见 \figref{rsa-flow}。

\shot[0.95]{fig-crypto-06-rsa-flow.png}{第 2 章 RSA 加解密流程：公钥 (n, e) 加密、私钥 d 解密，以及 d 与 φ(n) 的关系（原理示意图）}{rsa-flow}

\paragraph{一张图背后的四件事}
\begin{enumerate}
\item $n = p \times q$：模数是两个大素数的乘积。$p, q$ 是秘密；
\item 加密：$c \equiv m^{e} \pmod n$。任何人都能用公开的 $(n, e)$ 加密；
\item 解密：$m \equiv c^{d} \pmod n$。需要\kw{私钥指数 $d$}；
\item $d$ 怎么来：$\varphi(n) = (p-1)(q-1)$，$d$ 是 $e$ 模 $\varphi(n)$ 的逆元，即 $e \cdot d \equiv 1 \pmod{\varphi(n)}$。
\end{enumerate}
\paragraph{为什么 $d$ 能解密}$d$ 是 $e$ 模 $\varphi(n)$ 的\kw{逆元}，意思是 $e \cdot d \equiv 1 \pmod{\varphi(n)}$，也就是 $e \cdot d = 1 + k\,\varphi(n)$。解密就是把密文 $c = m^{e}$ 再取 $d$ 次幂，而指数相乘可以拆开：
\[
c^{d} = (m^{e})^{d} = m^{e \cdot d} = m^{1 + k\,\varphi(n)} = m \cdot \big(m^{\varphi(n)}\big)^{k} \equiv m \pmod n .
\]
这里用到欧拉定理 $m^{\varphi(n)} \equiv 1 \pmod n$（当 $\gcd(m,n)=1$），多出来的部分被“绕”成 $1^{k} = 1$，最后只剩 $m^{1} = m$。一句话直觉：\kw{逆元把乘法“倒回去”}——$e$ 把 $m$ 幂了一次，$d$ 正好把这次幂抵消，于是 $c^{d}$ 回到 $m$。前提仍是 $\gcd(e, \varphi(n)) = 1$，否则 $d$ 根本不存在。

\kw{私钥为什么不能公开}：知道 $d$ 就能解密一切。而要算 $d$ 就得知道 $\varphi(n)$，就得分解 $n$。所以“$n$ 很大且难分解”是 RSA 安全的根基。反过来，题目给的 $n$ 要么太小、要么被重复使用，根基就塌了——这正是 C3 和 C4 的突破口。

\paragraph{逐字节加密还是整体加密}要特别检查消息是怎样变成整数的。本地 C3 \kw{逐字节加密}：每个明文字节单独做一次模幂，所以每个密文整数对应一个字节，解出的整数应在 0--255。C4 则把\kw{整个字节串按大端序转换成一个大整数}再做一次模幂，解密后要用正确的字节序还原整个字节串。两者解密后的还原方式不同。Python 的 \code{int.from\_bytes} 与 \code{int.to\_bytes} 可以控制字节序（“大端”= 高位字节在前）；长度未知时须根据题面或文件格式推断，不能随便丢掉开头的零字节。

\memory{RSA 题先找 $(n, e, c)$ 三元组，再问两个问题：$n$ 能不能分解？$n$ 是不是被用了不止一次？}

\section{小模数为什么危险：试除分解}
当 $n$ 很小时（比如 $n = 3233$），分解它只需要\kw{试除}：从 2 开始逐个试 $n$ 能不能被整除，试到 $\lfloor\sqrt{n}\rfloor$ 就够了——因为 $n = p \times q$ 的两个因子不可能都大于 $\sqrt{n}$。找到 $p$ 就得到 $q = n / p$，接着算 $\varphi(n)$、求 $d$、解密。

$n = 3233$ 时 $\sqrt{n} \approx 56.8$，试到 56 找到 $53$，另一个因子是 $61$。这就是 C3 的全部数学。真正的 RSA 用几百位甚至几千位的 $n$，试除完全不可行；真实 RSA 还依赖正确填充等工程细节。\kw{不能把教学题的试除法当作普遍攻击}。

\pitfall{$\sqrt{n}$ 是试除的上限，不是随便试。超过 $\lfloor\sqrt{n}\rfloor$ 还没找到因子，说明 $n$ 是素数——这本身就是重要信息。}

\section{共模攻击为什么成立：同一 $n$、两个 $e$}
如果同一段明文 $m$ 被同一个模数 $n$ 下的两个不同指数 $e_1, e_2$ 加密，得到 $c_1 \equiv m^{e_1}$、$c_2 \equiv m^{e_2} \pmod n$，并且 $\gcd(e_1, e_2) = 1$，那么可以不分解 $n$ 就恢复 $m$。原理是数论里的\kw{扩展欧几里得算法}：能找到整数 $a, b$ 使
\[
e_1 a + e_2 b = 1.
\]
于是
\[
c_1^{a} \cdot c_2^{b} \equiv m^{e_1 a + e_2 b} \equiv m^{1} \equiv m \pmod n.
\]
负的系数（比如 $a = -8$）表示取\kw{模逆元}：$c_1^{-8}$ 先求 $c_1$ 的逆再乘 8 次幂。C4（玄机 587）正是这种情形：两条回执记录共用一个 308 位十进制的 $n$，指数分别是 $65537$ 与 $17$，且 $65537 \times (-8) + 17 \times 30841 = 1$。

\paragraph{那对系数 $a, b$ 是怎么算出来的}上式的 $a, b$ 由\kw{扩展欧几里得算法}给出，它是“辗转相除”的加强版，分两步理解：
\begin{enumerate}
\item \kw{普通辗转相除}求 $\gcd$：反复用“大数除以小数、取余数，再把除数当新的大数”，直到余数为 0，最后一个非零余数就是最大公约数。例如求 $\gcd(65537, 17)$：$65537 = 3855 \times 17 + 2$，$17 = 8 \times 2 + 1$，$2 = 2 \times 1 + 0$，所以 $\gcd(65537, 17) = 1$。
\item \kw{扩展版}在回代时顺带记录每一步，把余数写成两个原数的整数组合，最后得到 $x\,a + y\,b = \gcd(x, y)$——这条等式叫\kw{Bézout（贝祖）等式}。
\end{enumerate}
当 $\gcd(e_1, e_2) = 1$ 时，Bézout 等式右边正好是 1，代入前一个公式就把两个模幂拼成“指数为 1”的一次幂，直接得到 $m$，全程不碰 $n$ 的分解。\figref{eea-demo} 是真实运行：\code{math.gcd(12, 18) = 6}（不互素，逆元不存在）与 \code{math.gcd(65537, 17) = 1} 形成对照，辗转相除只走了 3 步，扩展版给出 $a = -8$、$b = 30841$，代回验证 $65537 \times (-8) + 17 \times 30841 = 1$。这也解释了为什么“指数互素”是共模攻击的前提：不互素时 Bézout 右边是 $\gcd > 1$，就拼不出指数为 1 的 $m$。

\shot[0.95]{fig-crypto-24-eea-demo.png}{第 2 章 gcd、辗转相除与扩展欧几里得：由 Bézout 系数 $a=-8$、$b=30841$ 得 $65537a+17b=1$，并演示逆元（真实终端实录）}{eea-demo}

\memory{同一明文 + 同一模数 + 两个互素指数 = 不用分解 $n$ 也能解密。}

\section{自定义符号编码：五格电文的方法论}
有些题不用标准编码，而是给一套自定义符号表。玄机 586“五格电文”的附件给出 32 个互不相同的符号、一份公开试报和一段待译正文。方法论只有一句话：\kw{先猜表示方式，再用已知样本验证}。
\begin{enumerate}
\item \kw{数一数符号种类}：32 个符号，而 $2^5 = 32$，最自然的假设是“每个符号表示 5 位”；
\item \kw{找已知明文}：附件给了公开试报，明文是“夜班OK”；
\item \kw{建立映射并验证}：每个符号换成它在字模板中的下标、写成 5 位二进制，从左到右每 8 位拼成一个字节。用试报做实验，能解出“夜班OK”，说明符号顺序、位序、字节编码都猜对了；
\item \kw{同法处理正文}：映射验证通过后，把完全相同的变换用于待译正文。
\end{enumerate}
注意 5 位符号的边界和 8 位字节的边界\kw{不对齐}，所以必须先把符号换成比特串、拼接、再按 8 位重新分组，见 \figref{fivegrid-bits}。抄报员加的空格只是分组记号，处理时一律去掉。

\shot[0.95]{fig-crypto-07-fivegrid-bits.png}{第 2 章 五格电文比特流示意：符号→下标→5 位→拼接→8 位一格→UTF-8 解码出“夜班OK”（原理示意图，数据来自 586 真实附件）}{fivegrid-bits}

\pitfall{看到 32 个符号不要直接套标准 Base32。标准 Base32 的字母表、填充规则都是固定的，自定义编码必须用自己的字模板和已知样本验证，不能假设“和 Base32 一样”。}

\section{本章知识点速查}
\begin{longtable}{p{0.28\linewidth}p{0.62\linewidth}}
\toprule
知识点 & 一句话要点 \\
\midrule
字符/字节/十六进制 & 字符给人看，字节给计算机存，十六进制给眼睛看字节；\code{b'41'} 有 2 字节，\code{fromhex('41')} 只有 1 字节。 \\
编码 vs 加密 & Base64、十六进制都是编码：规则公开、无密钥；加密有密钥。 \\
异或 & $(x \oplus k) \oplus k = x$；单字节密钥只有 256 种，可穷举。 \\
可读性判断 & 看前缀、字符范围、语言、结构四条一起看，别只看一个字符。 \\
模运算 & 除以 $n$ 取余；\code{pow(b, e, n)} 快速幂；\code{pow(a, -1, n)} 求逆元。 \\
RSA 三元组 & $(n, e, c)$ 公开；$d$ 私有，$ed \equiv 1 \pmod{\varphi(n)}$，$n = pq$。 \\
小模数 & $n$ 小到 $\lfloor\sqrt{n}\rfloor$ 能试完，试除即可分解。 \\
共模 & 同 $n$ 双 $e$ 且互素，用 Bézout 系数组合出明文。 \\
自定义编码 & 先猜表示方式，再用已知样本验证；5 位符号与 8 位字节边界不对齐。 \\
\bottomrule
\end{longtable}

% ===========================================================================
\chapter{Crypto 工具箱}
% ===========================================================================
\goal{三个主力工具上手：python 的一行命令与脚本、xxd 看十六进制、pow 三参数做快速幂与求逆元。}

本章每个工具都按“是什么 / 怎么装 / 基本用法 / 真实输出”介绍。工具不多，但每一个都会在第 4 章反复出现。

\section{python3：交互、一行命令与脚本}
\paragraph{是什么}Python 是本册唯一的“编程工具”。解题的三件事——读文件、做变换、打印核对量——都靠它。3.8 以上版本即可，标准库就够，不用装第三方包。

\paragraph{怎么装}Windows 在 \url{https://www.python.org/downloads/} 下载安装包，安装时勾选 \code{Add python.exe to PATH}。WSL 里执行 \code{sudo apt install python3}。装完用 \code{python --version} 验证。

\paragraph{基本用法}三种形态，按需要选：
\begin{itemize}
\item \kw{一行命令}：\code{python -c '...'} 适合两三行就说完的小实验；
\item \kw{脚本文件}：把代码存成 \code{xxx.py} 再 \code{python xxx.py}，适合超过三行、要反复修改的解题代码；
\item \kw{交互模式}：直接敲 \code{python} 进入 \code{>>>} 提示符，适合随手验证一个表达式（输 \code{exit()} 离开）。
\end{itemize}
\figref{tool-python} 是真实执行的三条命令：查版本、一行命令循环打印、导入标准库算最大公约数。

\shot[0.95]{fig-crypto-08-tool-python.png}{第 3 章 python 基本用法：版本确认、\code{-c} 一行命令、标准库随手用（真实终端实录）}{tool-python}

\subsubsection*{PowerShell 引号规则（看懂截图里的命令）}
截图里的一行命令写成 \texttt{python -c}\,\code{'print(''hello'')'}，为什么单引号都是两个？因为 PowerShell 的单引号字符串里，\kw{一个单引号要写两遍}才表示一个单引号字符；外层单引号把整段代码括成一个参数传给 python。所以 \code{''hello''} 到 python 手里就是 \code{'hello'}。照抄截图里的命令就能得到同样的输出。

\pitfall{一行命令超过三四行就写成脚本文件。一行命令里改一个逗号都容易改错，脚本文件可以用编辑器检查、反复运行。}

\section{xxd：看文件的十六进制视图}
\paragraph{是什么}\code{xxd} 把文件的原始字节按十六进制打印，右侧还带 ASCII 预览列。它是“看清文件到底存了什么字节”的最直接工具。在 WSL/Ubuntu 里自带（\code{sudo apt install xxd} 可补装），Windows 侧可用 WSL 运行。

\paragraph{基本用法}
\begin{lstlisting}
xxd 文件名          # 完整十六进制视图
xxd 文件名 | head   # 只看前若干行（WSL/bash）
\end{lstlisting}
\figref{tool-xxd} 是真实执行：对 \code{labs/crypto/xor.hex} 运行 \code{xxd}，注意左侧十六进制和右侧 ASCII 列——右侧显示的是 \code{515b5650...} 这些\kw{字符}，说明这个文件里存的是十六进制\kw{文本}，不是密文字节本身；对 \code{rsa.txt} 运行则看到 \code{n=3233...} 的明文参数。这一眼就能把“文件内容”和“文件显示”分开。

\shot[0.95]{fig-crypto-09-tool-xxd.png}{第 3 章 xxd 查看文件：xor.hex 存的是十六进制文本，rsa.txt 存的是可读参数（真实终端实录，WSL）}{tool-xxd}

\memory{\code{xxd} 右侧 ASCII 列有可读字符 = 这个文件存的是文本；全是一堆随机符号才是原始二进制字节。}

\section{pow 三参数：快速幂与求逆元}
\paragraph{是什么}Python 内置的 \code{pow(b, e, n)} 计算 $b^{e} \bmod n$，先做快速幂再取模，数值多大都不怕；\code{pow(a, -1, n)} 计算 $a$ 模 $n$ 的逆元（Python 3.8 起支持）。这两个是 RSA 题的日常操作。

\paragraph{基本用法}见 \figref{tool-pow} 的真实执行：
\begin{lstlisting}
pow(2, 10, 1000)      # 2^10 = 1024，模 1000 得 24
pow(17, -1, 3120)     # 17 模 3120 的逆元 = 2753
(17 * 2753) % 3120    # 验证乘回去余 1
math.gcd(65537, 17)   # 判断两个指数是否互素
\end{lstlisting}
注意 \code{pow(2, 10)} 两参数版本才是普通乘方（得 1024）；三参数版本才是模幂。RSA 解密就是一句 \code{pow(c, d, n)}。

\shot[0.95]{fig-crypto-10-tool-invmod.png}{第 3 章 \code{pow} 用法：快速幂、求逆元、验证逆元、判断互素（真实终端实录）}{tool-pow}

\paragraph{求逆元的前提}$a$ 与 $n$ 必须互素（最大公约数为 1），否则逆元不存在，\code{pow(a, -1, n)} 会报 \code{ValueError}。RSA 里要求 $\gcd(e, \varphi(n)) = 1$ 就是这个道理。

\section{读附件的三件小工具}
\paragraph{zipfile：读压缩包}平台题附件是 ZIP。不用解压也能读：
\begin{lstlisting}[language=Python]
from zipfile import ZipFile
z = ZipFile('labs/platform/attachments/586-five-grid.zip')
print(*z.namelist(), sep='\n')        # 列出包内文件
data = z.read('attachments/archive.json')  # 读出文件内容（字节）
\end{lstlisting}

\paragraph{json：读结构化数据}附件里的 \code{archive.json}、\code{export.json} 是 JSON 格式（键值对组成的文本）：
\begin{lstlisting}[language=Python]
import json
record = json.loads(data)        # 字节或文本 → Python 字典
print(record['glyph_plate'])     # 按键取值
\end{lstlisting}

\paragraph{ast：读 Python 风格的数据列表}\code{rsa.txt} 里的密文写成 \code{ciphertext=[1369, 745, ...]}，可以直接解析成列表：
\begin{lstlisting}[language=Python]
import ast
text = open('labs/crypto/rsa.txt').read()
data = ast.literal_eval(text.split('ciphertext=', 1)[1].strip())
\end{lstlisting}
\code{ast.literal\_eval} 只解析字面量（数字、列表、字符串），不会把数据当任意 Python 代码执行，比 \code{eval} 安全。

\section{工具箱小结}
\begin{longtable}{p{0.3\linewidth}p{0.3\linewidth}p{0.3\linewidth}}
\toprule
想做什么 & 用什么 & 关键点 \\
\midrule
跑一行小实验 & \code{python -c '...'} & PowerShell 里单引号写两遍表示一个单引号 \\
把字节显示成十六进制 & \code{xxd 文件}（WSL） & 看 ASCII 列判断文件存的是文本还是字节 \\
算 $b^{e} \bmod n$ & \code{pow(b, e, n)} & 快速幂，不怕大数 \\
求 $a^{-1} \bmod n$ & \code{pow(a, -1, n)} & 需要 $a$ 与 $n$ 互素 \\
十六进制转字节 & \code{bytes.fromhex(s)} & 长度必须是偶数 \\
字节转十六进制 & \code{b.hex()} & 1 字节变 2 个字符 \\
字节串转整数 & \code{int.from\_bytes(b, 'big')} & 记得字节序 \\
读 ZIP/JSON/listing & \code{zipfile} / \code{json} / \code{ast.literal\_eval} & 都是标准库 \\
\bottomrule
\end{longtable}

% ===========================================================================
% ===========================================================================
\chapter{题目：题面卡与逐题图解带做}

\goal{五道题，每题一节：先用题面卡把“给什么、要什么、考什么”读清楚，紧接着就是这道题的完整解题步骤——读题与思路、分步操作与观察、结果与核对、为什么会成功、排错指南、变式练习。建议先只读题面卡自己动手试 10--15 分钟，卡住再看后面的步骤。}

% ===========================================================================

每张卡的字段固定为：题目来源 / 给定材料 / 目标 / 交付物 / 考点 / 前置知识 / 难度。难度用星表示：★☆☆ 入门、★★☆ 进阶、★★★ 挑战。读完卡片先自己想两分钟，再往下看这道题的解题步骤。

% ===========================================================================

\paragraph{做题顺序建议}C0 $\rightarrow$ C1 $\rightarrow$ C3 $\rightarrow$ C2 $\rightarrow$ C4。C0、C1 练表示与穷举，C3 练 RSA 基本操作，C2 练自定义编码的验证方法，C4 练数论关系。卡住超过 20 分钟就看第 4 章对应小节的“读题与思路”。

% ===========================================================================
\section{题 C0：十六进制字符串到底有多少字节？}

\begin{challenge}{题 C0 题面卡}
\field{题目来源}{本地练习（离线可做），题面即本卡}
\field{给定材料}{一段文本 \code{666c61677b6865787d}}
\field{目标}{判断它的字符数、按十六进制还原后的字节数，并读出最终 ASCII 内容}
\field{交付物}{一个 \code{flag\{...\}} 字符串；提交前核对字符数与花括号}
\field{考点}{数据表示：十六进制是字节的显示方式，不是加密}
\field{前置知识}{2.1 节 字符/字节/十六进制；3.1 节 python 一行命令}
\field{难度}{★☆☆（入门）}
\end{challenge}

\paragraph{读题与思路}交付物是 flag 字符串，给的是一串十六进制字符。看到这种材料，第一反应不是“解密”，而是“换表示”：十六进制通常只是字节的显示方式。所以要问三个问题：现在手里的是文本还是字节？按两字符一字节还原后有多长？还原出的字节能不能按 ASCII 读出文字？题面还特意问“字符数、字节数、ASCII 内容”，这就是在提醒你别把字符数当字节数。

\paragraph{分步操作与观察}
\begin{enumerate}[label=\arabic*.]
\item 在 PowerShell 里让 python 把字符串按十六进制还原，并打印字符数与字节数。命令（单引号内的单引号写两遍，见 3.1 节）：
\begin{lstlisting}
python -c 's=''666c61677b6865787d''; print(''文本长度(字符):'', len(s));
b=bytes.fromhex(s); print(''还原后字节数:'', len(b)); print(''字节内容:'', b);
print(''ASCII 解码:'', b.decode(''ascii''))'
\end{lstlisting}
\figref{c0-hex} 是真实执行结果。应该看到：文本长度 18、还原后字节数 9、字节内容显示成 \code{b'flag\{hex\}'}。这一步是把“十六进制文本”换回“字节”。
\item 对比两种“41”，确认自己没有把显示方式当内容：
\begin{lstlisting}
python -c 'b1=b''41''; b2=bytes.fromhex(''41'');
print(''两个字符的文本 b41:'', b1, ''长度'', len(b1));
print(''解析成 1 个字节后:'', b2, ''长度'', len(b2))'
\end{lstlisting}
同一张 \figref{c0-hex} 的第二条命令就是它：\code{b'41'} 长度 2，\code{b'A'} 长度 1。应该看到两个长度不同。
\item 用 \code{.decode('ascii')} 把 9 个字节读成文本。ASCII 只定义了 0--127 的字符，flag 里的字母、花括号都在其中，能无损读出。
\end{enumerate}

\shot[0.95]{fig-crypto-11-c0-hex-bytes.png}{第 4 章 C0：18 个十六进制字符还原成 9 个字节，再按 ASCII 读出（真实终端实录）}{c0-hex}

\paragraph{结果与核对}最终 ASCII 内容是 \code{flag\{hex\}}。独立核对：把还原出的字节重新编码回十六进制，应得到原字符串；字符数 $18 = 9 \times 2$，每字节两个十六进制字符，数量关系对得上。

\flagbox{flag\{hex\}}

\paragraph{为什么会成功}一句话机制：\kw{十六进制是字节的显示方式，不是加密}。展开说：这道题没有任何密钥和算法，只有一次表示转换。18 个十六进制字符按“两个字符一组”解读，每组是 1 个字节的值；9 个字节按 ASCII 表解读成 9 个字符。理解了这层，就能理解为什么 Crypto 题的第一步永远是“识别表示”。

\debugtip{
\begin{itemize}
\item \code{bytes.fromhex} 报 \code{non-hexadecimal number}：字符串里混进了空格、换行或非十六进制字符，先 \code{.strip()} 并检查长度是否为偶数；
\item 长度是奇数：多复制或少复制了一个字符，回题面重新数；
\item 解码出乱码：确认用的是 \code{decode('ascii')} 或 \code{decode('utf-8')}，且数据没有被当作十六进制处理两次。
\end{itemize}}

\paragraph{变式练习}
\begin{enumerate}
\item 把 \code{666c6167} 还原成字节并读出文字，记录每一步的类型与长度（提示：4 个字符 → 4 个字节？先数清楚）。
\item 给字符串 \code{flag\{hex\}} 用 \code{.encode().hex()} 编码回十六进制，核对是否得到原题数据。
\item 找一段 Base64 文本（如 \code{ZmxhZ3toZXh9}）用 \code{base64.b64decode} 解码，记录“字符数 $\to$ 字节数”的关系，并对比本题的十六进制。
\end{enumerate}

% ===========================================================================
\section{题 C1：单字节异或如何从密文找回明文？}

\begin{challenge}{题 C1 题面卡}
\field{题目来源}{本地靶场 \code{labs/crypto/xor.hex}（离线可做）}
\field{给定材料}{密文文件 \code{labs/crypto/xor.hex}；题目说明：密钥只有 1 个字节，循环作用于每个明文字节}
\field{目标}{找出密钥，还原明文}
\field{交付物}{密钥（十六进制）与一个 \code{flag\{...\}} 字符串；两样都要写出核对过程}
\field{考点}{异或的可逆性；256 种密钥的穷举与可读性判断}
\field{前置知识}{2.2 节 异或真值表、2.3 节 单字节穷举；3.2 节 xxd（可选）}
\field{难度}{★☆☆（入门）}
\end{challenge}

\paragraph{读题与思路}题目说密钥只有 1 个字节、循环作用于每个明文字节。这句翻译成搜索量就是：\kw{256 种候选}。256 种完全试得起，所以本题的正路是穷举 + 可读性筛选，不是猜。动手前先看材料本身：\code{labs/crypto/xor.hex} 是什么表示？文件名已经提示是十六进制。用 \code{Get-Content} 看它的内容（或 \code{xxd} 看它的字节），确认之后再写代码。筛选真明文的线索是题面格式：flag 通常以 \code{flag\{} 开头且全是可打印字符。

\paragraph{分步操作与观察}
\begin{enumerate}[label=\arabic*.]
\item 先看密文文件内容，再转成字节：
\begin{lstlisting}
Get-Content labs/crypto/xor.hex
python -c 't=open(''labs/crypto/xor.hex'').read().strip();
print(''十六进制字符数:'', len(t)); b=bytes.fromhex(t);
print(''字节数:'', len(b)); print(''前 8 字节的十六进制:'', b[:8].hex());
print(''直接打印字节串(乱码正常):'', b)'
\end{lstlisting}
\figref{c1-file} 是真实执行结果。应该看到：44 个十六进制字符、22 个字节。直接打印字节串会看到 \code{Q[VPLOXEh\textasciicircum DYXChZVP\textasciicircum TJ} 之类的乱码——这正是异或后的密文“不可读”的直观样子。（注意 \code{.strip()} 是为了去掉文件末尾的换行，不加会多出一个字节。）
\item 穷举 256 个密钥并观察候选：
\begin{lstlisting}
python -c 'data=bytes.fromhex(open(''labs/crypto/xor.hex'').read().strip());
plain=lambda k: bytes(x^k for x in data);
[print(hex(k), ''->'', plain(k).decode(''ascii'', errors=''replace''))
 for k in range(0x33, 0x3c)];
hit=[k for k in range(256)
     if plain(k).startswith(b''flag{'') and all(32<=c<127 for c in plain(k))];
print(''--- 256 个候选全部试完, 全可读且以 flag{ 开头的密钥:'', [hex(k) for k in hit]);
print(''密钥'', hex(hit[0]), ''明文'', plain(hit[0]).decode())'
\end{lstlisting}
\figref{c1-brute} 是真实执行结果。打印 \code{0x33} 到 \code{0x3b} 这一段可以看到“差一点”的候选全是乱码，只有 \code{0x37} 那一行是 \code{flag\{xor\_is\_not\_magic\}}；最后的筛选行显示 256 个候选中只有 \code{'0x37'} 满足“全部可打印 + 以 \code{flag\{} 开头”。这一步是穷举与可读性判断。
\item 用两个独立方法核对密钥（\figref{c1-verify}）：
\begin{lstlisting}
python -c 'data=bytes.fromhex(open(''labs/crypto/xor.hex'').read().strip());
k=data[0]^ord(''f''); print(''由首字节 f 反推密钥:'', hex(k));
plain=bytes(x^k for x in data); print(''明文:'', plain.decode());
back=bytes(x^k for x in plain);
print(''明文再异或得到的十六进制:'', back.hex());
print(''与原密文一致:'', back.hex()==open(''labs/crypto/xor.hex'').read().strip())'
\end{lstlisting}
应该看到：首字节法也得到 \code{0x37}；明文再异或一遍得到的十六进制与原密文逐字符一致，打印 \code{True}。
\end{enumerate}

\shot[0.95]{fig-crypto-12-c1-xor-hex.png}{第 4 章 C1 第 1 步：xor.hex 是十六进制文本，44 个字符换 22 个字节（真实终端实录）}{c1-file}

\shot[0.95]{fig-crypto-13-c1-brute.png}{第 4 章 C1 第 2 步：256 个单字节密钥穷举，只有 0x37 命中可读 flag（真实终端实录）}{c1-brute}

\shot[0.95]{fig-crypto-14-c1-verify.png}{第 4 章 C1 第 3 步：首字节反推密钥 0x37，反向异或与原密文一致（真实终端实录）}{c1-verify}

\paragraph{结果与核对}密钥是 \code{0x37}（十进制 55），明文是 \code{flag\{xor\_is\_not\_magic\}}。三重独立核对：穷举筛选唯一命中、首字节公式 $k = c_0 \oplus \mathrm{ord}(f)$ 同样得到 0x37、明文反向异或还原出原密文。

\flagbox{flag\{xor\_is\_not\_magic\}}

\paragraph{为什么会成功}一句话机制：\kw{$(x \oplus k) \oplus k = x$，且单字节密钥空间只有 256}。展开说：异或加密把明文每个字节和同一个 $k$ 做异或；解密再做一次就还原。因为 $k$ 只有 256 种取值，攻击者可以全部试一遍，用“可打印 + \code{flag\{} 前缀”把假候选排掉。首字节法是同一原理的捷径：已知明文第一个字节是 \code{f}，就能直接算出 $k$。但请注意：这种方法只在\kw{密钥空间小或重复使用}时成立；若密钥与消息等长、真正随机、只用一次，穷举就没有意义。

\debugtip{
\begin{itemize}
\item \code{bytes.fromhex} 报奇数长度：多半是读文件时带了换行，先 \code{.strip()}；
\item 穷举结果全是乱码、没有一个像 flag：确认异或对象是“字节”而不是十六进制“文本”；别对 \code{t} 直接异或，要先 \code{bytes.fromhex(t)}；
\item 筛选出多个候选：加强筛选条件，检查完整 \code{flag\{} 前缀、花括号配对和字符范围，短数据尤其容易误判；
\item \code{decode('ascii')} 报错：加 \code{errors='replace'} 先看内容，或改用 \code{errors='ignore'} 并记录丢掉了什么。
\end{itemize}}

\paragraph{变式练习}
\begin{enumerate}
\item 自己写一段 \code{flag\{my\_name\}} 的明文，用密钥 \code{0x42} 异或加密，把密文十六进制交给同学（或隔一天的自己）用穷举解回。
\item 把题目改成 3 字节循环密钥：位置 $0,3,6,\dots$ 共用 $k_0$，位置 $1,4,7,\dots$ 共用 $k_1$，位置 $2,5,8,\dots$ 共用 $k_2$。提示：按下标对 3 取余分组，每组就是一道单字节异或题。解完用反向异或核对整段。
\item 只用首字节法解本题（$k = c_0 \oplus \mathrm{ord}(\mathrm{f})$），不跑穷举。再故意假设首字节是 \code{F}（大写）重算，观察结果为何不可信，体会“已知明文假设也要验证”。
\end{enumerate}

% ===========================================================================
\section{题 C2：五格电文的符号到底表示什么？（玄机 586）}

\begin{challenge}{题 C2 题面卡}
\field{题目来源}{玄机平台 \emph{2026安网杯-五格电文}，\url{https://xj.edisec.net/challenges/586}；附件本地副本 \url{labs/platform/attachments/586-five-grid.zip}（离线可做到解出 flag）}
\field{给定材料}{附件 ZIP 内有字模板 \code{glyph\_plate}（32 个互不相同的符号）、公开试报 \code{sample\_wire} 与已知明文 \code{sample\_plain}（“夜班OK”）、待译正文 \code{dispatch\_wire}，以及说明文件 \code{format.txt}（抄报员以空格断句，誊录时一律省略）}
\field{目标}{推出符号到比特的映射，翻译正文，读出其中的 flag}
\field{交付物}{一个 \code{flag\{...\}} 字符串（32 位十六进制样式），在玄机平台对应题目页提交}
\field{考点}{自定义符号编码：先猜表示方式（32 符号 $\approx$ 5 位），再用已知样本验证映射与位序}
\field{前置知识}{2.1 节 编码与表示；2.8 节 自定义编码方法论；3.4 节 zipfile/json}
\field{难度}{★★☆（进阶）}
\end{challenge}

\paragraph{读题与思路}材料是 32 个怪符号组成的电文、一份公开试报和一段待译正文。看到“32 个符号”，第一个假设应该是\kw{每个符号表示 5 位}，因为 $2^5 = 32$。但注意：这只是假设。验证假设的钥匙在附件里——公开试报的明文已知（“夜班OK”），它是你的\kw{标尺}。所以正确的动作顺序是：数符号 $\rightarrow$ 提假设 $\rightarrow$ 用试报验证 $\rightarrow$ 验证通过才解正文。反过来“看到 32 个字符就套标准 Base32”是错的：标准 Base32 的字母表是固定的 \code{A-Z2-7}，而这里的符号是自定义的，必须用字模板自己建立映射。

\paragraph{分步操作与观察}
\begin{enumerate}[label=\arabic*.]
\item 打开附件副本，先看包里有什么：
\begin{lstlisting}
python -c 'from zipfile import ZipFile;
z=ZipFile(''labs/platform/attachments/586-five-grid.zip'');
print(*z.namelist(), sep=chr(10))'
\end{lstlisting}
\figref{c2-zip} 上半部分是真实执行结果：包里有 \code{attachments/archive.json}（数据）和 \code{attachments/format.txt}（说明）。说明文件写明：字模板上的符号次序与电文取值相关；抄报员以空格断句，仅为换气记号，誊录时一律省略。
\item 读出数据并核对符号表的前提条件：
\begin{lstlisting}
python -c 'from zipfile import ZipFile; import json;
r=json.loads(ZipFile(''labs/platform/attachments/586-five-grid.zip'')
              .read(''attachments/archive.json''));
print(''archive.json 的字段:'', sorted(r));
print(''glyph_plate:'', r[''glyph_plate'']);
print(''符号数:'', len(r[''glyph_plate'']), '', 去重后:'', len(set(r[''glyph_plate''])));
print(''sample_wire:'', r[''sample_wire'']);
print(''dispatch_wire 前 80 字符:'', r[''dispatch_wire''][:80])'
\end{lstlisting}
\figref{c2-zip} 下半部分是真实执行结果。应该看到：字段有 \code{dispatch\_wire}、\code{glyph\_plate}、\code{sample\_plain}、\code{sample\_wire}、\code{station}；符号数 32、去重后还是 32。\kw{互不重复}很关键——它保证“符号 $\to$ 下标”的映射是唯一确定的。
\item 用公开试报验证映射（本题的核心一步）：
\begin{lstlisting}
python -c 'from zipfile import ZipFile; import json;
r=json.loads(ZipFile(''labs/platform/attachments/586-five-grid.zip'')
              .read(''attachments/archive.json''));
a=r[''glyph_plate'']; w=r[''sample_wire''].replace('' '', '''');
bits=''''.join(f''{a.index(s):05b}'' for s in w);
print(''试报符号数:'', len(w), '', 比特数:'', len(bits));
print(''前 12 个符号在模板中的下标:'', [a.index(s) for s in w[:12]]);
print(''拼出的比特串前 40 位:'', bits[:40]);
b=bytes(int(bits[i:i+8], 2) for i in range(0, len(bits)//8*8, 8));
print(''每 8 位组装成字节:'', b); txt=b.decode(''utf-8'');
print(''UTF-8 解码结果:'', txt);
print(''与 sample_plain 一致:'', txt==r[''sample_plain''])'
\end{lstlisting}
\figref{c2-sample} 是真实执行结果。注意看四件事：去掉空格后 13 个符号 $\to$ 65 个比特；每个符号换成字模板下标再写成 5 位二进制；每 8 位拼成一个字节（\code{b'\textbackslash xe5\textbackslash xa4\textbackslash x9c...'}）；UTF-8 解码得到 \code{夜班OK}，与 \code{sample\_plain} 比对打印 \code{True}。映射、位序、字节编码三项同时验证通过。
\item 同一个变换处理正文并输出 flag：
\begin{lstlisting}
python labs/platform/solve_five_grid.py
python labs/platform/solve_five_grid.py labs/platform/attachments/586-five-grid.zip
\end{lstlisting}
\figref{c2-run} 是真实执行结果（两条命令分别是“默认读本地附件”和“显式给 ZIP 路径”）。脚本会先自动核对公开试报再输出正文，你应该看到 \code{sample: 夜班OK}、完整正文和 \code{flag: flag\{c619...\}}。
\item 到玄机平台题目页 \url{https://xj.edisec.net/challenges/586} 提交 flag，等待提交按钮结束加载，确认页面显示“FLAG 正确”。
\end{enumerate}

\shot[0.95]{fig-crypto-15-c2-zip-contents.png}{第 4 章 C2 第 1--2 步：586 附件结构、字模板 32 个符号互不重复、试报与正文（真实终端实录）}{c2-zip}

\shot[0.95]{fig-crypto-16-c2-sample-check.png}{第 4 章 C2 第 3 步：公开试报解出“夜班OK”，与 sample\_plain 一致 True（真实终端实录）}{c2-sample}

\shot[0.95]{fig-crypto-17-c2-five-grid-run.png}{第 4 章 C2 第 4 步：solve\_five\_grid.py 输出正文与 flag（真实终端实录）}{c2-run}

\paragraph{结果与核对}正文末尾的通行标记是：
\flagbox{flag\{c619020e518c599051efc804f5f84884\}}
三重核对：公开试报解出“夜班OK”验证了映射；正文里 flag 符合 32 位十六进制样式；本账号此前在平台提交后显示正确反馈。样本报错时脚本会直接抛异常，不会输出正文——这就是“验证不过就不解正文”的自动化体现。

\paragraph{为什么会成功}一句话机制：\kw{已知明文是验证自定义编码假设的标尺}。展开说：32 个符号给出“每符号 5 位”的候选假设；字模板给出“符号 $\to$ 数值”的映射；公开试报给出正确答案。把三者拼起来就能验证“下标顺序 + 位序 + 每 8 位分组 + UTF-8”整条链路，任何一环错了都会在试报上露馅。链路验证通过后，正文只是同一条流水线的输入而已。另外注意 5 位边界和 8 位边界不对齐（\figref{fivegrid-bits}），所以必须先拼比特流再按 8 位重新分组，不能“两个符号凑一个字节”。

\debugtip{
\begin{itemize}
\item 试报解出来不对：先检查符号表顺序（是不是照抄字模板原文、有没有去重或排序）、再检查位序（下标写 5 位时高位在前）、最后检查是否去掉了所有空格；
\item 解出的字节无法 UTF-8 解码：分组边界错了。打印比特串长度，确认按 $8k$ 位截断，余数位应丢弃（\code{len(bits)//8*8}）；
\item 脚本报 \code{sample does not match}：映射没验证通过，正文先别解，回头查试报；
\item 脚本报 \code{glyph\_plate must contain 32 distinct symbols}：符号表读错了字段或数据损坏，重新解压附件。
\end{itemize}}

\paragraph{变式练习}
\begin{enumerate}
\item 把字模板倒序（或循环左移 3 位）生成一份新题，先故意用错映射解试报，观察报错；再改回正确映射解出正文。体会“样本验证”抓错的效率。
\item 不看脚本，手算 \code{sample\_wire} 的前 3 个符号：写下它们在字模板中的下标、5 位二进制，和截图里的比特串前 40 位逐段对照。
\item 思考题：如果题目不给公开试报，只给正文，你还能验证映射吗？列出至少两个可以利用的弱假设（比如正文通常是 UTF-8 中文、结尾有 \code{flag\{...\}} 样式），并说明它们为什么只是弱假设。
\end{enumerate}

% ===========================================================================
\section{题 C3：小模数 RSA 为什么能被试除？}

\begin{challenge}{题 C3 题面卡}
\field{题目来源}{本地靶场 \code{labs/crypto/rsa.txt}（离线可做）}
\field{给定材料}{参数文件 \code{labs/crypto/rsa.txt}：$n = 3233$、$e = 17$、逐字节密文整数列表（23 项）}
\field{目标}{恢复明文，并说清这道题弱在何处}
\field{交付物}{一个 \code{flag\{...\}} 字符串 + 一段“弱点说明”（哪个数学条件没守住）}
\field{考点}{RSA 三元组；小模数试除分解；求逆元与逐字节解密}
\field{前置知识}{2.4 节 模运算、2.5 节 RSA 三元组、2.6 节 小模数；3.3 节 \code{pow}}
\field{难度}{★★☆（进阶）}
\end{challenge}

\paragraph{读题与思路}材料是三个数 $n = 3233$、$e = 17$ 和 23 个密文整数。看到 RSA 先找三元组 $(n, e, c)$，然后问两个问题（2.5 节的口诀）：$n$ 能不能分解？$n$ 被用了几次？这里 $n = 3233$ 只有四位数，$\sqrt{n} \approx 56.8$，试除最多试 55 次就能分解——\kw{可以分解}。再看题面说“逐字节密文列表”，意味着每个明文字节单独加密一次，解密出的每个整数应该都落在 0--255。所以路线是：分解 $n$ $\to$ 算 $\varphi(n)$ $\to$ 求 $d$ $\to$ 逐项 \code{pow(c, d, n)}。

\paragraph{分步操作与观察}
\begin{enumerate}[label=\arabic*.]
\item 读入参数并检查密文列表的合法性：
\begin{lstlisting}
Get-Content labs/crypto/rsa.txt
python -c 'import ast;
t=open(''labs/crypto/rsa.txt'').read();
data=ast.literal_eval(t.split(''ciphertext='', 1)[1].strip());
print(''密文整数个数:'', len(data));
print(''最大值:'', max(data), '', 是否都小于 n=3233:'', all(c<3233 for c in data));
print(''前 6 项:'', data[:6])'
\end{lstlisting}
\figref{c3-file} 是真实执行结果。应该看到 23 个整数、最大值 2923、\code{True}。每项都小于 $n$ 是应该的：模幂的输出本就在 $0$ 到 $n-1$ 之间。
\item 试除分解 $n$ 并求私钥 $d$：
\begin{lstlisting}
python -c 'import math; n=3233;
fs=[p for p in range(2, math.isqrt(n)+1) if n%p==0];
print(''从 2 试除到'', math.isqrt(n), '', 找到因子:'', fs, '', 另一因子:'', [n//p for p in fs]);
phi=(61-1)*(53-1); d=pow(17, -1, phi);
print(''phi(n) ='', phi, '', 私钥指数 d ='', d);
print(''验证 (17*d) % phi ='', (17*d) % phi)'
\end{lstlisting}
\figref{c3-trial} 上半部分是真实执行结果：试除到 56 找到因子 $53$，另一因子 $61$；$\varphi(n) = 3120$；$d = 2753$；验证 $(17 \times 2753) \bmod 3120 = 1$。这一步把“分解 $\to$ 欧拉函数 $\to$ 逆元”三件事串成一条链，每件都当场核对。
\item 用往返测试确认公式没错，再解密：
\begin{lstlisting}
python -c 'n=3233; c=pow(65, 17, n);
print(''演示 pow(65,17,3233) ='', c, '', 再 pow(c,2753,3233) ='', pow(c, 2753, n))'
\end{lstlisting}
同图 \figref{c3-trial} 下半部分：随便取一个明文数 65，加密再解密回到 65，说明 $(e, d)$ 配对正确。
\begin{lstlisting}
python -c 'import ast;
t=open(''labs/crypto/rsa.txt'').read();
data=ast.literal_eval(t.split(''ciphertext='', 1)[1].strip());
pt=bytes(pow(c, 2753, 3233) for c in data);
print(''解密得到的前 6 个整数:'', [pow(c, 2753, 3233) for c in data[:6]]);
print(''明文:'', pt.decode());
back=[pow(b, 17, 3233) for b in pt];
print(''用公钥再加密与原密文一致:'', back==data)'
\end{lstlisting}
\figref{c3-decrypt} 是真实执行结果。前 6 个整数是 $102, 108, 97, 103, 123, 114$，正是 \code{f}、\code{l}、\code{a}、\code{g}、\code{\{}、\code{r} 的 ASCII 码；整段解码得到 flag；再加密比对打印 \code{True}。
\end{enumerate}

\shot[0.95]{fig-crypto-18-c3-rsa-txt.png}{第 4 章 C3 第 1 步：rsa.txt 的参数与 23 项密文整数，全部小于 n（真实终端实录）}{c3-file}

\shot[0.95]{fig-crypto-19-c3-pow-trialdiv.png}{第 4 章 C3 第 2--3 步：试除分解 3233=61×53、求 d=2753、往返测试（真实终端实录）}{c3-trial}

\shot[0.95]{fig-crypto-20-c3-decrypt.png}{第 4 章 C3 第 4 步：逐字节解密得到 flag，公钥再加密核对 True（真实终端实录）}{c3-decrypt}

\paragraph{结果与核对}明文是：
\flagbox{flag\{rsa\_small\_modulus\}}
核对有三层：$17 \times 2753 \bmod 3120 = 1$ 确认 $d$ 是 $e$ 的逆元；往返测试 $65 \to 2790 \to 65$ 确认公式；整段明文用公钥再加密得到的列表与 \code{rsa.txt} 逐项一致（打印 \code{True}）。\kw{弱点说明}：这道题弱在模数太小，试除就能分解出 $p, q$；真实 RSA 还会保证参数足够大并使用填充方案，不能把本题的试除法当作普遍攻击。

\paragraph{为什么会成功}一句话机制：\kw{$n$ 小到能分解，RSA 的安全根基就没了}。展开说：RSA 的安全性建立在“分解 $n$ 很难”之上；$d$ 的计算需要 $\varphi(n) = (p-1)(q-1)$，而分解出 $p, q$ 就能算出 $\varphi(n)$。试除的复杂度是 $O(\sqrt{n})$，$n = 3233$ 时只有几十次；换成 1024 位的 $n$，试除次数是天文数字。另外注意“逐字节加密”这个出题设定：它让每个密文整数只对应一个字节，解密结果天然落在 0--255，验证非常容易。

\debugtip{
\begin{itemize}
\item 解出的整数有大于 255 的：说明“逐字节”假设或 $d$ 算错了。先核对 $17d \bmod \varphi(n) = 1$，再核对 $pq = n$；
\item 解出来是乱码：检查是否用了正确的密文列表（别把 $n$、$e$ 也读进去）、是否每项单独解密、是否整体转成大整数解密后再按字节序还原（本题不是这种）；
\item \code{pow(17, -1, phi)} 报 \code{ValueError}：$e$ 与 $\varphi(n)$ 不互素，说明分解或 $\varphi$ 算错了；
\item 列表解析报错：\code{rsa.txt} 里 \code{ciphertext=} 后面才是列表，用 \code{split('ciphertext=', 1)[1]} 切开再解析。
\end{itemize}}

\paragraph{变式练习}
\begin{enumerate}
\item 自己选两个小素数（比如 $p = 101, q = 113$）生成一组新的 $(n, e)$ 和逐字节密文，写成同款 \code{rsa.txt}，用本题流程解回来。
\item 把素数逐步增大（$101 \to 1009 \to 10007$），用 \code{time} 模块测每次试除分解耗时，观察成本随 $\sqrt{n}$ 上升的趋势。
\item 把逐字节加密改成“整段字节串按大端序转成一个大整数再模幂”，对比两种做法解密后的还原步骤有什么不同（提示：\code{int.to\_bytes} 要指定长度和字节序）。
\end{enumerate}

% ===========================================================================
\section{题 C4：同一模数的两份回执能否互相补全？（玄机 587）}

\begin{challenge}{题 C4 题面卡}
\field{题目来源}{玄机平台 \emph{2026安网杯-联号回执}，\url{https://xj.edisec.net/challenges/587}；附件本地副本 \url{labs/platform/attachments/587-linked-receipts.zip}（离线可做到解出 flag）}
\field{给定材料}{附件 ZIP 内 \code{export.json} 有两条记录：同一回执正文的两份 RSA 密文，\code{modulus} 相同（308 位十进制），\code{exponent} 分别是 $65537$ 与 $17$；说明文件写明正文按 UTF-8 编码后转成大端无符号整数再模幂}
\field{目标}{不分解大模数，利用两份密文之间的关系恢复正文}
\field{交付物}{一个 \code{flag\{...\}} 字符串（32 位十六进制样式），在玄机平台对应题目页提交}
\field{考点}{RSA 共模攻击：同 $n$ 双互素指数 $\Rightarrow$ 扩展欧几里得组合出明文；负指数与模逆元}
\field{前置知识}{2.5 节 RSA 三元组、2.7 节 共模攻击；3.3 节 \code{pow}；3.4 节 zipfile/json}
\field{难度}{★★★（挑战）}
\end{challenge}

\paragraph{读题与思路}材料是两条回执记录：同一个模数 $n$、指数分别是 $65537$ 和 $17$、两份不同的密文。看到“同一明文、同一模数、两个指数”，应该立刻想到\kw{共模}（2.7 节）；看到大模数（308 位十进制），应该立刻放弃分解。动手前先核对三个前提：$n_1 = n_2$ 吗？$\gcd(e_1, e_2) = 1$ 吗？明文确实是同一份吗（说明文件写明“两条记录保护同一份回执正文”）？前提都成立，才能用 Bézout 系数组合。

\paragraph{分步操作与观察}
\begin{enumerate}[label=\arabic*.]
\item 读附件，核对两条记录的公钥参数（\figref{c4-fields}）：
\begin{lstlisting}
python -c 'from zipfile import ZipFile; import json;
recs=json.loads(ZipFile(''labs/platform/attachments/587-linked-receipts.zip'')
                 .read(''attachments/export.json''))[''records''];
print(''记录数:'', len(recs));
[print(r[''receipt_id''], '', e ='', r[''exponent''], '', n 的十进制位数 ='',
  len(r[''modulus'']), '', 密文位数 ='', len(r[''ciphertext''])) for r in recs];
print(''两条记录模数相同:'', recs[0][''modulus'']==recs[1][''modulus'']);
import math;
print(''两个指数互素:'', math.gcd(int(recs[0][''exponent'']), int(recs[1][''exponent'']))==1)'
\end{lstlisting}
应该看到：两条记录 RS-2048-A（$e = 65537$）与 RS-2048-B（$e = 17$），$n$ 都是 308 位十进制数，模数相同打印 \code{True}，指数互素打印 \code{True}。
\item 求 Bézout 系数：找 $a, b$ 使 $65537a + 17b = 1$（\figref{c4-bezout}）：
\begin{lstlisting}
python -c 'import sys; sys.path.insert(0, ''labs/platform'');
from solve_receipt import bezout;
g, a, b = bezout(65537, 17);
print(''gcd(65537, 17) ='', g, '', a ='', a, '', b ='', b);
print(''验证 65537*a + 17*b ='', 65537*a + 17*b)'
\end{lstlisting}
真实输出：\code{a = -8}、\code{b = 30841}，代回验证得 1。负系数无妨：$c_1^{-8}$ 表示先求 $c_1$ 的模逆元再乘 8 次幂，Python 的 \code{pow(c1, -8, n)} 直接支持。
\item 运行复现脚本解出正文（\figref{c4-run}）：
\begin{lstlisting}
python labs/platform/solve_receipt.py
python labs/platform/solve_receipt.py labs/platform/attachments/587-linked-receipts.zip
\end{lstlisting}
应该看到三行输出：\code{coefficients: -8 30841}、完整正文、\code{flag: flag\{f6d5...\}}。脚本在输出前会做两件事：检查前提（模数相同、指数互素、密文与 $n$ 互素）和\kw{再加密核对}（把恢复的 $m$ 分别用两个指数重新加密，必须等于附件里的两份密文），任一步不过都会抛异常而不是打印乱码。
\item 到玄机平台题目页 \url{https://xj.edisec.net/challenges/587} 提交 flag，确认页面显示“FLAG 正确”。
\end{enumerate}

\shot[0.95]{fig-crypto-21-c4-receipt-fields.png}{第 4 章 C4 第 1 步：587 两条记录同模数、指数 65537 与 17 互素（真实终端实录）}{c4-fields}

\shot[0.95]{fig-crypto-22-c4-bezout.png}{第 4 章 C4 第 2 步：扩展欧几里得给出 a=-8、b=30841，负指数用模逆元处理（真实终端实录）}{c4-bezout}

\shot[0.95]{fig-crypto-23-c4-receipt-run.png}{第 4 章 C4 第 3 步：solve\_receipt.py 输出正文与 flag（真实终端实录）}{c4-run}

\paragraph{结果与核对}正文末尾的通行标记是：
\flagbox{flag\{f6d5f18a35d8aa5802b12f7312b18b63\}}
核对链：$\gcd(65537, 17) = 1$；$65537 \times (-8) + 17 \times 30841 = 1$；恢复的 $m$ 用两个指数再加密后与附件密文逐项一致（脚本内建检查）；本账号此前在平台提交后显示正确反馈。

\paragraph{为什么会成功}一句话机制：\kw{指数互素时，两个模幂可以组合成一次“指数为 1”的模幂}。展开说：$c_1 = m^{65537}$、$c_2 = m^{17} \pmod n$，而 $65537a + 17b = 1$，所以 $c_1^a c_2^b = m^{65537a + 17b} = m$。攻击全程没有分解 $n$，突破口是\kw{参数复用}而不是大数分解。这与 C3 的“小模数可分解”是两种完全不同的弱点——看到 RSA 字样不能跳过附件条件检查直接套用某一种攻击。

\debugtip{
\begin{itemize}
\item 脚本报 \code{the two records have different moduli}：两条记录不是共模，共模推导不成立，回题面核对附件版本；
\item 脚本报 \code{common-modulus conditions are not met}：指数不互素，或某份密文与 $n$ 不互素（$\gcd(c, n) \neq 1$ 时该密文根本没被正确生成）；
\item 脚本报 \code{re-encryption check failed}：算出的 $m$ 不对，检查 Bézout 系数方向（$a$ 配 $e_1$、$b$ 配 $e_2$）和负指数的处理；
\item 解出乱码或缺字节：整段字节串按大端还原时长度算错。用 \code{to\_bytes((m.bit\_length()+7)//8, 'big')} 这种按位长度换算，不要随手丢零字节。
\end{itemize}}

\paragraph{变式练习}
\begin{enumerate}
\item 构造小规模共模练习：取 $n = 3233$、明文一个字节，分别用 $e_1 = 3$、$e_2 = 5$ 加密，用扩展欧几里得找系数后解回明文，并做再加密核对。
\item 把两指数换成不互素的组合（如 $e_1 = 3$、$e_2 = 9$）重做上题，观察推导在哪一步断掉，并解释为什么“不互素”意味着两个方程提供的信息不足以恢复 $m$。
\item Bézout 系数不唯一：从 $a = -8, b = 30841$ 出发，构造另一组解（提示：$a' = a + 17t,\ b' = b - 65537t$），用它再解一次本题，验证结果相同。
\end{enumerate}

% ===========================================================================

% ===========================================================================

\chapter{复盘：知识地图与易错清单}
% ===========================================================================
\goal{把五道题串成一张“下次遇到这类题怎么办”的地图，并把最容易踩的坑列成清单。}

\section{解题路线图：拿到材料先做什么}
五道题的完整决策流程如下，建议自己默画一遍：
\begin{enumerate}
\item \kw{第 1 步：识别表示}。现在手里的是什么？
\begin{itemize}
\item 只有 \code{0-9a-fA-F} 且长度偶数 $\rightarrow$ 先 \code{bytes.fromhex}（C0）；
\item 文件里是十六进制文本 $\rightarrow$ 先 \code{Get-Content} / \code{xxd} 看真实内容（C1）；
\item 一堆怪符号、种类有限 $\rightarrow$ 数种类，猜“每符号几位”（C2：32 种 $\rightarrow$ 5 位）；
\item 若干个十进制大整数 + $n, e$ $\rightarrow$ RSA 三元组（C3、C4）。
\end{itemize}
\item \kw{第 2 步：找可逆关系或弱点}。
\begin{itemize}
\item 密钥空间小 $\rightarrow$ 穷举（C1：单字节 256 种）；
\item $n$ 小 $\rightarrow$ 试除分解（C3：$3233 = 61 \times 53$）；
\item 同 $n$ 双 $e$ 且互素 $\rightarrow$ 共模组合（C4：$65537 \times (-8) + 17 \times 30841 = 1$）；
\item 自定义编码 $\rightarrow$ 用已知样本标定映射（C2：试报“夜班OK”）。
\end{itemize}
\item \kw{第 3 步：验证}。
\begin{itemize}
\item 单字节异或 $\rightarrow$ 反向异或还原原密文；
\item RSA $\rightarrow$ 用公钥再加密比对密文；
\item 自定义编码 $\rightarrow$ 用已知明文验证映射；
\item 平台题 $\rightarrow$ 提交后看“FLAG 正确”。
\end{itemize}
\item \kw{第 4 步：写题解}。保留输入版本、关键常数来源、算法、运行命令和平台反馈；把失败过的假设也写下来。
\end{enumerate}

\section{下次遇到这类题怎么办}
\begin{longtable}{p{0.3\linewidth}p{0.6\linewidth}}
\toprule
题型信号 & 固定动作 \\
\midrule
一串十六进制字符 & 数字符数 $\to$ \code{bytes.fromhex} $\to$ 数字节数 $\to$ 尝试 ASCII/UTF-8 解码。记录每一步长度。 \\
一段不可读的字节、题面说“单字节密钥” & \code{bytes.fromhex} 后遍历 0--255，按“\code{flag\{} 前缀 + 全可打印”筛选；再用反向异或核对。 \\
有限种符号组成的电文 & 数符号种类猜位数（$2^5 = 32$）；找已知明文样本先验证映射；符号转比特、拼接、按 8 位分组。 \\
$(n, e, c)$ 且 $n$ 很小 & 试除到 $\lfloor\sqrt{n}\rfloor$ 分解 $\to$ $\varphi(n)$ $\to$ \code{pow(e, -1, phi)} $\to$ \code{pow(c, d, n)}。 \\
同一明文、同一 $n$、两个互素 $e$ & 扩展欧几里得找 $a, b$，用 $c_1^a c_2^b \bmod n$ 组合；负指数走模逆元；最后再加密核对。 \\
结果拿不准是不是 flag & 按四条查：前缀、字符范围、结构（花括号配对）、独立再验证。 \\
\bottomrule
\end{longtable}

\section{易错清单（按“错一下浪费半小时”的顺序）}
\begin{enumerate}
\item \kw{把十六进制文本当字节}：\code{b'41'} 是 2 字节，\code{bytes.fromhex('41')} 才是 1 字节。动手前先数长度。
\item \kw{读文件忘 \code{.strip()}}：文件末尾的换行也是一个字节，会让异或和长度全错。
\item \kw{对文本而不是字节做异或}：先把十六进制解析成字节，再异或。
\item \kw{只检查一个字符就下结论}：短数据有很多“像真的”候选，要查完整前缀和整段字符范围。
\item \kw{忘了 $\lfloor\sqrt{n}\rfloor$ 是试除上限}：试到 56 为止（$n = 3233$），不是无限试。
\item \kw{$\varphi(n)$ 算错}：$\varphi(n) = (p-1)(q-1)$，不是 $pq - 1$。算完先验证 $ed \bmod \varphi(n) = 1$。
\item \kw{看到 RSA 直接上某一种攻击}：先检查附件条件——$n$ 小不小？$n$ 用没用过两次？指数互不互素？不同条件用不同攻击。
\item \kw{共模攻击忘核对前提}：$n_1 = n_2$、$\gcd(e_1, e_2) = 1$、明文相同，缺一不可。
\item \kw{负指数不会处理}：$c^{-8} \bmod n$ 直接写 \code{pow(c, -8, n)}（Python 3.8+），或先 \code{pow(c, -1, n)} 再乘幂。
\item \kw{字节序还原错误}：整体转整数的 RSA 要用 \code{int.to\_bytes} 按大端还原，长度按位数换算，别乱丢零字节。
\item \kw{验证环节省略}：脚本打印 flag 不等于对；再加密/反向运算/已知样本三类核对至少做一类，平台题还要看提交反馈。
\item \kw{把猜测写成答案}：没被验证的假设写进题解时必须标注“未验证”，这是复盘能复用的前提。
\end{enumerate}

\memory{先识别表示，再找可逆关系，最后独立验证。三步都做完，答案才可信。}

% ===========================================================================
\chapter{自测与参考答案}
% ===========================================================================
\goal{先做自测题再看答案。答案里的推演刻意写全中间判断，直接抄最后一行是自欺。}

\section{课后自测题}
\begin{enumerate}
\item \code{bytes.fromhex('666c6167')} 有几个字节？把它按 ASCII 读出来是什么？
\item 为什么说 Base64 是编码不是加密？举一个“编码没有密钥”的判断依据。
\item 写出异或的四行真值表，并用一行公式说明“加密 = 解密”。
\item 重复异或密钥长度为 3 时，密文的哪些位置共享同一个密钥字节？如何把这类题拆成子题？
\item 为什么单字节异或可以穷举？“可读性判断”至少要检查哪四条？
\item 用一句话定义模运算，并写出 Python 中计算 $b^{e} \bmod n$ 与 $a^{-1} \bmod n$ 的表达式。
\item RSA 的公钥三元组是什么？私钥指数 $d$ 与 $\varphi(n)$ 是什么关系？
\item 什么条件下 $e$ 在模 $\varphi(n)$ 下有逆元？如果不满足会怎样？
\item 为什么 $n = 3233$ 能被试除分解？试除的上限为什么是 $\lfloor\sqrt{n}\rfloor$？
\item 逐字节 RSA 与“整体转换成整数”的 RSA，在还原明文时有什么区别？
\item 共模攻击成立需要哪三个条件？写出组合明文的公式。
\item 五格电文为什么要先解公开试报？如果直接套标准 Base32 会发生什么？
\end{enumerate}

\section{参考答案与完整推演}

\paragraph{第 1 题}4 个字节（8 个十六进制字符 $\div$ 2），读出来是 \code{flag}。每个字节 2 个十六进制字符，所以字符数永远是字节数的两倍。

\paragraph{第 2 题}Base64 的转换规则完全公开，任何人不需要秘密信息就能解码；加密必须有密钥，没有密钥就还原不了。判断依据就是“是否需要秘密密钥”。同理十六进制、URL 编码都是编码。

\paragraph{第 3 题}
\begin{lstlisting}
x  k  | x^k
0  0  | 0
0  1  | 1
1  0  | 1
1  1  | 0
\end{lstlisting}
一行公式：$(x \oplus k) \oplus k = x$。异或自身可逆，所以加密程序即解密程序。

\paragraph{第 4 题}位置 $0, 3, 6, 9, \dots$ 共用 $k_0$；位置 $1, 4, 7, \dots$ 共用 $k_1$；位置 $2, 5, 8, \dots$ 共用 $k_2$。按“下标 mod 3”把密文分成三组，每组就是一道单字节异或题，分别穷举后按原位置拼回，最后用反向异或核对整段。

\paragraph{第 5 题}因为密钥只有 1 个字节，取值 $0$--$255$ 共 256 种，可以全部试完。四条检查：完整的 \code{flag\{} 前缀；字符是否全在可打印范围（32--126）；语言规律（英文单词/UTF-8 中文合法性）；结构合理性（花括号配对、长度合适）。

\paragraph{第 6 题}模运算 = 除以 $n$ 取余数，结果在 $0$ 到 $n-1$ 之间。Python：\code{pow(b, e, n)} 与 \code{pow(a, -1, n)}（要求 $\gcd(a, n) = 1$）。

\paragraph{第 7 题}三元组是 $(n, e, c)$：模数、加密指数、密文。$d$ 是 $e$ 模 $\varphi(n)$ 的逆元，即 $e \cdot d \equiv 1 \pmod{\varphi(n)}$；$\varphi(n) = (p-1)(q-1)$，其中 $n = pq$。

\paragraph{第 8 题}当且仅当 $\gcd(e, \varphi(n)) = 1$ 时 $e$ 有模 $\varphi(n)$ 逆元。不满足则 $d$ 不存在，标准 RSA 解密公式失效（出题人若故意这样设置，通常另有蹊跷，如 $e$ 与 $\varphi(n)$ 共享因子时的特例分析，超出本册范围）。

\paragraph{第 9 题}$n = 3233$ 只有四位数，$\sqrt{3233} \approx 56.8$，最多试 55 个候选即可。上限是 $\lfloor\sqrt{n}\rfloor$ 因为 $n = p \times q$ 的两个因子不可能同时大于 $\sqrt{n}$：若 $p > \sqrt{n}$ 且 $q > \sqrt{n}$，则 $pq > n$，矛盾。

\paragraph{第 10 题}逐字节 RSA：每个密文整数独立对应一个明文字节，解密后直接 \code{bytes(...)} 组装，每个结果必须在 0--255。整体转整数的 RSA：解密得到一个大整数，要用 \code{int.to\_bytes(length, 'big')} 按字节序还原整个字节串，长度要按位数换算、零字节不能乱丢。

\paragraph{第 11 题}三条件：同一明文；同一模数 $n$；$\gcd(e_1, e_2) = 1$。公式：先求 $a, b$ 使 $e_1 a + e_2 b = 1$（扩展欧几里得），则 $m \equiv c_1^{a} \cdot c_2^{b} \pmod n$，负指数用模逆元。

\paragraph{第 12 题}公开试报是唯一能\kw{验证假设}的材料：它同时校验符号表顺序、位序、分组方式和字符编码。直接套标准 Base32 会错在两处：字母表不对（这里的符号是自定义的 \code{glyph\_plate}），填充规则未知。任何一步不匹配都应退回检查假设，而不是硬解正文。

\section{完整推演：C3 的每一步为何必要}
C3 附件给出 $n = 3233$、$e = 17$ 和 23 个密文整数。先确认列表中每项都小于 $n$，因为模幂的输出本应在 $0$ 到 $n-1$ 之间；这一步是合法性检查，防止把别的数据误当密文。再分解 $n$：从 2 试除到 $\lfloor\sqrt{3233}\rfloor = 56$，找到 $53$，于是 $3233 = 61 \times 53$。因此 $\varphi(n) = 60 \times 52 = 3120$。求解 $17d \equiv 1 \pmod{3120}$，得 $d = 2753$；核对 $17 \times 2753 = 46801 = 15 \times 3120 + 1$，对 3120 取余确为 1。公式链闭合。

题目生成脚本提示“逐字节”，于是每个密文整数都对应一个明文字节。脚本中使用 \code{pow(c, d, n)} 计算模幂，比先计算巨大的 $c^{d}$ 再取余高效且清晰。把所有结果放进 \code{bytes(...)} 后按 ASCII/UTF-8 解码，得到可读 flag。若某个结果超过 255，说明假设或参数有误；若得到乱码，先检查是否应整体解密、是否用错字节序、是否读错密文列表，而不是任意删字符。
\begin{lstlisting}[language=Python]
from math import isqrt
n = 3233
factors = [(p, n // p) for p in range(2, isqrt(n) + 1)
           if n % p == 0]
print(factors)                      # [(53, 61)]
phi = (61 - 1) * (53 - 1)           # 3120
d = pow(17, -1, phi)                # 2753
print(d, (17 * d) % phi)            # 2753 1
\end{lstlisting}
正常 RSA 不应选这样容易试除的模数。本题的目的在于把公钥、分解、模逆与解密串成可验证的链条，每一步都有独立核对点：$pq = n$、$ed \bmod \varphi(n) = 1$、往返测试、再加密比对。

\section{完整推演：C4 的推导为何成立}
587 的两条记录共享模数 $n$，指数 $e_1 = 65537$、$e_2 = 17$。先核对 $\gcd(65537, 17) = 1$，于是扩展欧几里得算法给出 $65537 \times (-8) + 17 \times 30841 = 1$。由 $c_1 \equiv m^{65537}$、$c_2 \equiv m^{17} \pmod n$ 可得
\[
c_1^{-8} \cdot c_2^{30841} \equiv m^{65537 \times (-8) + 17 \times 30841} \equiv m \pmod n.
\]
负指数表示模逆元：$c_1^{-8} = (c_1^{-1})^{8}$，计算前要确认 $c_1$ 与 $n$ 互素。得到 $m$ 后按大端转换回字节并用 UTF-8 解码。最后重新计算两次加密，检查它们是否分别等于附件中的密文，避免“乱码里偶然出现 flag 样式文本”的误判。附带脚本会执行条件检查、解密和再加密核对：
\begin{lstlisting}
python labs/platform/solve_receipt.py
\end{lstlisting}
\textbf{答案：}\code{flag\{f6d5f18a35d8aa5802b12f7312b18b63\}}，已在平台提交并获得正确反馈。这个案例与 C3 的“小模数可分解”突破口完全不同：只看到 RSA 字样，不能跳过附件条件检查直接套用一种攻击。

\section{完整推演：C2 的验证链为何必要}
586 的字模板有 32 个互不重复的符号，“每符号 5 位”成为可验证假设的依据是 $2^5 = 32$。去掉 \code{sample\_wire} 中的分组空格，把每个符号替换为它在模板中的下标、用五位二进制写出，再从左到右每八位还原一个字节。这个方法解出 \code{sample\_plain} 中的“夜班OK”，说明符号顺序、位序与 UTF-8 解码均与已知样本相符；然后用完全相同的函数处理 \code{dispatch\_wire}，得到完整正文，末尾包含所需 flag。
\begin{lstlisting}
python labs/platform/solve_five_grid.py
\end{lstlisting}
\textbf{答案：}\code{flag\{c619020e518c599051efc804f5f84884\}}，已在平台提交并获得正确反馈。关键不是“看到 32 个字符就套 Base32”，而是用公开试报核对符号表、位序和字节编码；任何一步不匹配都应退回去检查假设。

% ===========================================================================
\chapter*{附录 A：玄机平台真题状态与提交 checklist}
\addcontentsline{toc}{chapter}{附录 A：玄机平台真题状态与提交 checklist}
% ===========================================================================

\section{本册涉及的平台题状态}
\begin{longtable}{p{0.2\linewidth}p{0.24\linewidth}p{0.46\linewidth}}
\toprule
题号与名称 & 本地材料 & 状态（照实记录，绝不编造） \\
\midrule
586 五格电文 & 附件 ZIP、\code{solve\_five\_grid.py} & 断网可解出并打印完整 flag；\textbf{已在平台提交并获正确反馈}。 \\
587 联号回执 & 附件 ZIP、\code{solve\_receipt.py} & 断网可解出并打印完整 flag；\textbf{已在平台提交并获正确反馈}。 \\
\bottomrule
\end{longtable}
说明：平台题面、价格与界面可能变化，始终以当前题目详情页为准；本讲义不保存账号口令。提交后必须看到页面显示“FLAG 正确”及步骤完成，才能把答案标记为“已验证”。未验证的推测不得写成答案。

\section{提交 checklist}
\begin{enumerate}[label=\arabic*.]
\item 题目与链接是否记清：586 五格电文 \url{https://xj.edisec.net/challenges/586}，587 联号回执 \url{https://xj.edisec.net/challenges/587}；附件是否为当前官方版本，本地副本路径是否为 \url{labs/platform/attachments/586-five-grid.zip} 与 \url{labs/platform/attachments/587-linked-receipts.zip}。
\item 数学与结构前提是否核对：586 字模板 32 符号互不重复、试报解出“夜班OK”；587 两条记录 \code{modulus} 相同、指数 $65537$ 与 $17$ 互素、B\'ezout 系数满足 $65537(-8)+17(30841)=1$。
\item 是否做再加密/反向核对：587 对恢复的 $m$ 分别用两个指数重新加密，与附件两项密文逐项比较；586 用公开试报验证映射；C1 把明文重新异或；C3 把明文用公钥再加密。
\item 运行命令是否可移植：\texttt{python}\,\url{labs/platform/solve_five_grid.py}、\texttt{python}\,\url{labs/platform/solve_receipt.py}（也接受 ZIP 路径参数）；不在讲义与脚本中写账号口令，不打包复制平台附件给队友，只给链接与运行命令，由队友自行下载。
\item 平台反馈是否记录：提交后页面显示“FLAG 正确”及步骤完成；未验证的推测不写成答案，题解中保留已完成阶段与未完成步骤。
\end{enumerate}

% ===========================================================================
\chapter*{附录 B：命令速查表}
\addcontentsline{toc}{chapter}{附录 B：命令速查表}
% ===========================================================================
两列表：左列“想做什么”，右列“敲什么命令”（Windows PowerShell；标注 WSL 的在 Ubuntu 里执行）。

\begin{longtable}{p{0.42\linewidth}p{0.5\linewidth}}
\toprule
想做什么 & 敲什么命令 \\
\midrule
查 Python 版本 & \code{python --version} \\
跑一行 Python & \code{python -c 'print(2+3)'}（PowerShell 里单引号写两遍表示一个单引号） \\
跑脚本 & \code{python labs/platform/solve\_five\_grid.py} \\
进交互模式 & \code{python}（退出：\code{exit()}） \\
看文件内容 & \code{Get-Content labs/crypto/xor.hex}（WSL：\code{cat labs/crypto/xor.hex}） \\
看文件字节（十六进制） & \code{xxd labs/crypto/xor.hex}（WSL） \\
十六进制转字节 & \code{bytes.fromhex('515b5650')} \\
字节转十六进制 & \code{b.hex()} \\
字符转字节 & \code{'A'.encode('utf-8')} \\
字节转字符 & \code{b'A'.decode('ascii')} \\
数字符数/字节数 & \code{len(s)} / \code{len(b)} \\
单字节异或 & \code{bytes(x \textasciicircum{} k for x in data)} \\
算 $b^{e} \bmod n$ & \code{pow(2, 10, 1000)} \\
求逆元 $a^{-1} \bmod n$ & \code{pow(17, -1, 3120)} \\
判断互素 & \code{math.gcd(65537, 17)} \\
试除分解小 $n$ & \code{[p for p in range(2, math.isqrt(n)+1) if n\%p==0]} \\
整数转字节串 & \code{m.to\_bytes((m.bit\_length()+7)//8, 'big')} \\
字节串转整数 & \code{int.from\_bytes(b, 'big')} \\
列 ZIP 内容 & \code{ZipFile(p).namelist()} \\
读 JSON 附件 & \code{json.loads(z.read(p))} \\
解析 listing 形式数据 & \code{ast.literal\_eval(text)} \\
切到资料根目录 & \code{cd C:/Users/你的用户名/Desktop/CTF零基础自学资料包} \\
WSL 里进资料目录 & \code{cd /mnt/c/Users/你的用户名/Desktop/CTF零基础自学资料包} \\
\bottomrule
\end{longtable}

% ===========================================================================
\chapter*{附录 C：术语表}
\addcontentsline{toc}{chapter}{附录 C：术语表}
% ===========================================================================
按拼音排序，每个术语一句话解释；页码见目录对应章节。

\begin{longtable}{p{0.24\linewidth}p{0.68\linewidth}}
\toprule
术语 & 一句话解释 \\
\midrule
ASCII & 把 0--127 的数字分配给英文字母、数字、标点的老编码表，\code{A} 是 65。 \\
Base64/Base32 & 把字节换成可打印字符的\textbf{编码}，规则公开无密钥，不是加密。 \\
Bézout 系数 & 使 $e_1 a + e_2 b = \gcd(e_1, e_2)$ 的整数对 $(a, b)$，共模攻击的钥匙。 \\
flag & CTF 题的答案，通常形如 \code{flag\{...\}}。 \\
哈希摘要 & 把任意数据映射成固定长度串的单向函数输出，如 32 位十六进制的 MD5。 \\
可读性判断 & 用前缀、字符范围、语言、结构四条线索从候选中挑出真明文。 \\
模运算（mod） & 除以 $n$ 取余数，结果在 $0$ 到 $n-1$ 之间。 \\
模逆元 & 满足 $a \cdot a^{-1} \equiv 1 \pmod n$ 的 $a^{-1}$，要求 $a$ 与 $n$ 互素。 \\
穷举（brute force） & 把所有可能的密钥全试一遍；只在密钥空间小时可行。 \\
三元组 $(n, e, c)$ & RSA 题最常见的公开材料：模数、加密指数、密文。 \\
素数（质数） & 只能被 1 和自身整除的大于 1 的整数，RSA 用两个素数相乘造模数。 \\
试除 & 从 2 逐个试除到 $\lfloor\sqrt{n}\rfloor$ 的分解方法，只对小 $n$ 有效。 \\
私钥指数 $d$ & RSA 解密用的秘密指数，$e \cdot d \equiv 1 \pmod{\varphi(n)}$。 \\
位序 & 多位二进制中高低位的排列顺序，自定义编码常在这里设坑。 \\
小端序/大端序 & 多字节数据的存放顺序：低位在前叫小端，高位在前叫大端。 \\
一次性密码本 & 与消息等长、真随机、只用一次的密钥，理论上不可破。 \\
异或（XOR） & 按位“不同得 1”的运算，$(x \oplus k) \oplus k = x$。 \\
字节 & 计算机存储的基本单位，一个 0--255 的数。 \\
字节序 & 同“小端序/大端序”。 \\
自定义编码 & 出题人自定符号表与位宽的表示方式，必须用已知样本验证映射。 \\
\bottomrule
\end{longtable}

\paragraph{写在最后}做 Crypto 题的三个动作请练到成为习惯：\kw{识别表示}、\kw{找可逆关系}、\kw{独立验证}。下次拿到一串看不懂的数据，先别猜，先数长度、看字节、找结构。祝你解题顺利。

\end{document}
